\documentclass[11pt,a4paper]{article}
\usepackage[T1]{fontenc}
\usepackage{lmodern}
\usepackage[margin=27mm,headheight=15pt]{geometry}
\usepackage{amsmath,amssymb,amsthm,mathtools}
\usepackage{booktabs,longtable,array,enumitem,graphicx,xcolor,microtype}
\usepackage{xurl,fancyhdr}
\definecolor{navy}{RGB}{24,55,81}
\definecolor{teal}{RGB}{0,104,112}
\usepackage[colorlinks=true,linkcolor=navy,citecolor=teal,urlcolor=teal]{hyperref}
\usepackage{bookmark}
\hypersetup{pdftitle={From Polylogarithm Conjectures to Exact Reductions},
pdfauthor={Research continuation prepared for Vladimir Reshetnikov},
pdfsubject={Gaussian multiple polylogarithms, exact shuffle certificates, unit ladders, and log-gamma asymptotics}}
\pagestyle{fancy}
\fancyhf{}
\fancyhead[L]{\small\color{navy}Polylogarithms: exact reductions and asymptotics}
\fancyhead[R]{\small 9 October 2026}
\fancyfoot[C]{\thepage}
\renewcommand{\headrulewidth}{0.3pt}
\setlength{\parskip}{0.3em}
\setlength{\parindent}{1.25em}
\setlength{\emergencystretch}{3em}
\setcounter{tocdepth}{2}
\numberwithin{equation}{section}
\theoremstyle{plain}
\newtheorem{theorem}{Theorem}[section]
\newtheorem{lemma}[theorem]{Lemma}
\newtheorem{proposition}[theorem]{Proposition}
\newtheorem{corollary}[theorem]{Corollary}
\theoremstyle{definition}
\newtheorem{definition}[theorem]{Definition}
\newtheorem{conjecture}[theorem]{Conjecture}
\theoremstyle{remark}
\newtheorem{remark}[theorem]{Remark}
\newtheorem{example}[theorem]{Example}
\DeclareMathOperator{\Li}{Li}
\DeclareMathOperator{\Rea}{Re}
\DeclareMathOperator{\Ima}{Im}
\DeclareMathOperator{\Var}{Var}
\DeclareMathOperator{\Res}{Res}
\DeclareMathOperator{\rank}{rank}
\newcommand{\Q}{\mathbb Q}
\newcommand{\C}{\mathbb C}
\newcommand{\Z}{\mathbb Z}
\newcommand{\sh}{\mathbin{\shuffle}}
\newcommand{\shuffle}{\mathbin{\amalg}}
\newcommand{\ii}{\mathrm i}
\newcommand{\one}{\boldsymbol 1}
\newcommand{\LG}{\mathcal L}
\newcommand{\dd}{\,\mathrm d}
\newcommand{\code}[1]{\nolinkurl{#1}}

\begin{document}
\begin{center}
{\LARGE\bfseries\color{navy}From Polylogarithm Conjectures\\[3pt]
to Exact Reductions}\\[9pt]
{\large Shuffle certificates, complementary-power ladders,\\
and effective log-gamma moment asymptotics}\\[12pt]
{\normalsize Research continuation for the ProveIt project}\\[3pt]
{\normalsize Prepared for Vladimir Reshetnikov}\\[3pt]
{\normalsize 9 October 2026}
\end{center}

\begin{abstract}
We continue the arithmetic-bridge programme of the ProveIt polylogarithm
manuscript with explicit analytic proofs and reproducible finite
certificates. Eleven identities that the audited manuscript presents as
numerically supported receive local proofs: its weight-five Gaussian
double relation, five weight-six Gaussian double evaluations, three
weight-four Gaussian triple evaluations, and the cubic and quartic
algebraic-base trilogarithm ladders. The weight-five relation is a
two-product shuffle consequence; retaining both products improves the
same-argument spanning bound from three generators to two. An explicit
binomial transform proves the appropriate parity reductions at every
double weight, while a one-zero iterated-integral formula reduces all
indices consisting of one two and any number of ones.

We identify the exact formal shuffle quotient at every weight and depth
using its multigraded Lyndon count, and give an executable rational
normal-form procedure. The cubic ladder receives a certificate comprising
two Kummer substitutions, one Landen relation, and four Rogers pentagons;
the quartic ladder extends to every complementary pair of powers.
In the log-gamma programme we correct the literature status of the cubic
moment, give a rational interval certificate for 112 decimal places,
and study the coefficients of the known factorial moment expansion.
A Lagrange-inverse and local-limit argument determines their exponential
growth and \(k^{-3/2}\) factor, proves divergence at every fixed moment,
and yields an explicit superexponentially small remainder when the
truncation order grows with the moment. Classical tools and previously
known general theorems are attributed throughout. Formal quotient
dimensions and numerical checks are never used as arithmetic
independence proofs.
\end{abstract}

\noindent\textbf{Source checkpoint.}
The starting manuscript is
\href{https://github.com/VladimirReshetnikov/ProveIt/tree/afed07429d3d37eceb6c8e9e54cf4da2d3f39d53/Analysis/Polylogarithms/docs/manuscript}
{the ProveIt polylogarithms manuscript},
at commit \code{afed07429d3d37eceb6c8e9e54cf4da2d3f39d53}.
The companion archive records source hashes, exact certificates,
numerical diagnostics, and reproduction commands.

\newpage
{\small\setlength{\parskip}{0pt}\tableofcontents}
\newpage

\section{Scope, conventions, and the status of the results}
\label{sec:scope}

The starting point is the unified reading manuscript
\emph{Polylogarithms and their Arithmetic Bridges} \cite{manuscript}.
Its purpose is broader than the evaluation of isolated constants:
classical polylogarithms, colored nested sums, gamma functions,
Hurwitz-zeta jets, and algebraic ladders are different coordinates on
related analytic constructions. The present continuation selects three
places where those coordinates can be turned into short, checkable
proofs: the same-argument Gaussian depth problem, algebraic unit ladders,
and log-gamma moments.

The mathematical task is to explain why an identity holds. A small
decimal residual is useful for detecting a sign or transcription error,
but it cannot replace a derivation. Accordingly, every proved identity
below either follows from a displayed analytic argument or is accompanied
by exact rational operations on explicitly stated functional equations.
The remaining conjectures are kept separate from these results.

\subsection{What changes in the manuscript}

Table~\ref{tab:status} lists the precise changes of mathematical status.
The reference to a chapter is to the pinned source manuscript, not to
the section numbering of this article.

\begin{table}[htbp]
\centering\small
\begin{tabular}{>{\raggedright\arraybackslash}p{0.22\linewidth}>{\raggedright\arraybackslash}p{0.28\linewidth}>{\raggedright\arraybackslash}p{0.39\linewidth}}
\toprule
Source location & Previously recorded status & Result established here\\
\midrule
Chapter 4, weight-five Gaussian relation
& Numerical relation among four same-argument doubles
& Proof from two shuffle products; a spanning set of at most two quantities\\
\addlinespace
Chapter 4, five weight-six doubles
& Numerical evaluations
& All five proved by one explicit all-weight parity formula\\
\addlinespace
Chapter 4, three weight-four triples
& Numerical evaluations at \(i\)
& All three proved; an arbitrary-depth classical-reduction family\\
\addlinespace
Chapter 4, weight-six triple quotient
& Finite rank seven on ten coordinates
& Exact all-weight, all-depth formal theorem; the three stated coordinates are Lyndon generators\\
\addlinespace
Chapter 3, cubic and quartic bases
& Two trilogarithm formulas with numerical evidence
& Explicit functional-equation certificates for both\\
\addlinespace
Chapter 7, cubic log-gamma moment
& A named depth-two coordinate still sought
& Published Tornheim-derivative evaluation identified; certified independent numerical enclosure\\
\addlinespace
Log-gamma moment continuation
& Fixed-order factorial asymptotic expansion known in the literature
& Coefficient growth, divergence, and a remainder bound for a growing truncation order\\
\bottomrule
\end{tabular}
\caption{Exact scope of this continuation. None of the spanning statements
asserts numerical linear independence.}
\label{tab:status}
\end{table}

The eleven promoted identities are counted as one weight-five relation,
five weight-six relations, three triple relations, and two algebraic-base
relations. The extra weight-five shuffle equation is a further proved
identity and is what sharpens the spanning statement.
We do not count a repeated specialization of a general formula as a
separate solution of an open problem.

\subsection{Normalization and convergence}

We use decreasing summation indices:
\[
\Li_{s_1,\ldots,s_d}(z_1,\ldots,z_d)
=\sum_{n_1>\cdots>n_d>0}
\frac{z_1^{n_1}\cdots z_d^{n_d}}
{n_1^{s_1}\cdots n_d^{s_d}}.
\]
For a single varying argument, write
\[
F_{s_1,\ldots,s_d}(z)
=\Li_{s_1,\ldots,s_d}(z,1,\ldots,1).
\]
Where no confusion is possible we suppress the \(F\) and the trailing
ones. The weight is \(w=s_1+\cdots+s_d\), and the series depth is \(d\).
For \(|z|<1\), the series defining this one-variable family converges
absolutely. At a unit-circle point other than \(1\), the boundary values
used below agree with the ordinary convergent series; the needed
conditional double cases are justified in the parity proof. At higher
depth the inner multiple harmonic sum is \(O(\log^{d-1}n)\); its positive
successive increments, together with the factor \(n^{-s_1}\), give a
sequence tending to zero with finite total variation. Summation by parts
then proves convergence away from \(z=1\).

The forms
\[
\omega_0=\frac{\dd t}{t},\qquad
\omega_1=\frac{\dd t}{1-t}
\]
give the outer-letter-first word
\[
F_{s_1,\ldots,s_d}(z)
=I_{0^{s_1-1}1\cdots0^{s_d-1}1}(z).
\]
Here \(I_{a_1\cdots a_w}(z)=\int_0^z\omega_{a_1}(t)
I_{a_2\cdots a_w}(t)\), and \(I_{\varnothing}=1\).
Every displayed word ends in \(1\), so the origin is integrable.
For example, \(I_{01}=\Li_2\),
\(I_{011}=F_{2,1}\), and \(I_{101}=F_{1,2}\).
These examples fix the convention before any shuffle is used.

At the Gaussian point define
\[
g_{a,b}=\Im F_{a,b}(i),\quad
h_{a,b}=\Re F_{a,b}(i),\quad
G=\beta(2),\quad L=\log2,\quad
\lambda_n=\Im\Li_n((1+i)/2).
\]
The Dirichlet beta function is
\(\beta(s)=\sum_{n\ge0}(-1)^n/(2n+1)^s\).
All complex logarithms in the explicit half-point calculation are
principal values, with their continuation paths stated.
Logarithms of positive algebraic numbers in the ladder section are real.

\subsection{Prior results and the novelty boundary}

The general parity theorem is due to Panzer \cite{Panzer}, whose paper
already gives explicit formulas at depths two and three.
Umezawa \cite{Umezawa} subsequently supplied an explicit general-depth
formula, including regularized boundary formulations.
Our binomial expression is an independently proved specialization suited
to the manuscript's convention and to small exact certificates.
It is not a new general parity theorem.

The polynomial structure of a shuffle algebra on Lyndon generators is
classical \cite{Radford}; the rank formula below is its precise
application to the relevant one-variable subalgebra.
The Landen, Rogers, and Kummer equations used for the ladders are likewise
classical \cite{Zagier,NakamuraShiraishi}. What is supplied here is a
complete local route from those equations to the exact coefficients
recorded in the manuscript.

For log-gamma moments, the fixed-order factorial expansion was proved by
Amdeberhan et al.\ \cite{ACEMKM}, and the cubic Tornheim representation
was obtained by Bailey, Borwein, and Borwein \cite{BBB}.
The further coefficient asymptotic and the explicit growing-order
remainder estimate are developed here from those starting points.
These statements are presented with proofs and a limited novelty claim:
they extend the audited manuscript, while no assertion of exhaustive
priority over the entire inverse-gamma literature is made.

\section{Exact shuffle quotients and a two-product proof at weight five}
\label{sec:shuffle}

\subsection{The quotient being computed}

Let \(\mathcal H^1\) be the rational vector space on the empty word
and all binary words ending in \(1\), with shuffle multiplication.
It is bigraded by word length \(w\) and number of ones \(d\).
Let \(\mathcal J\) be its augmentation ideal, spanned by nonempty
words. The indecomposable quotient
\[
Q_{w,d}=(\mathcal J/\mathcal J^2)_{w,d}
\]
sets products of two positive-weight elements equal to zero.
This is a formal combinatorial construction. Shuffle equalities evaluate
to analytic identities at \(z\); the corresponding map on indecomposable
quotients requires also quotienting the evaluated algebra by products.
Evaluation can introduce further relations. In particular, the formal
dimension is not the dimension of a space of numerical periods.

\begin{theorem}[The exact formal shuffle count]\label{thm:witt}
For \(w\ge1\) and \(1\le d\le w\),
\[
\dim_{\Q}Q_{w,d}
=\ell_{w,d}
:=\frac1w\sum_{k\mid\gcd(w,d)}
\mu(k)\binom{w/k}{d/k}.
\]
A basis is represented by the binary Lyndon words of length \(w\)
with \(d\) ones. Every such word belongs to \(\mathcal H^1\).
Every other word has a finite, exact polynomial expression in the full
family of Lyndon generators, each monomial having the original total
bidegree. The number of independent product relations in
this component is
\[
\binom{w-1}{d-1}-\ell_{w,d}.
\]
\end{theorem}

\begin{proof}
Order the alphabet by \(0<1\) and words lexicographically.
A Lyndon word is strictly smaller than every nontrivial rotation.
Every word has a unique Chen--Fox--Lyndon factorization into a
nonincreasing concatenation of Lyndon words.
In characteristic zero, shuffle monomials in Lyndon words form a
polynomial basis \cite{Radford}. One can see the constructive point
directly: if the factorization is \(u_1^{m_1}\cdots u_r^{m_r}\) with
distinct \(u_1>\cdots>u_r\), then the shuffle product
\(u_1^{\shuffle m_1}\shuffle\cdots\shuffle u_r^{\shuffle m_r}\)
has this concatenation as its lexicographically largest word, with
coefficient \(m_1!\cdots m_r!\). All other words are smaller.
Induction therefore gives a triangular polynomial expression, and the
same triangularity proves linear independence of the monomials.

Every Lyndon word of length greater than one that contains both letters
starts with \(0\) and ends with \(1\). The single-letter words are
\(0\) and \(1\). The factorization of a word ending in \(1\) cannot
contain the factor \(0\), because that smallest factor would be last,
forcing the concatenation to end in \(0\).
Thus \(\mathcal H^1\) is exactly the polynomial algebra on all Lyndon
words other than \(0\). Passing to its augmentation quotient leaves
one class for each generator.

It remains to count. A primitive binary word of length \(w\) with
\(d\) ones has \(w\) distinct rotations and exactly one Lyndon rotation.
Every binary word is a power of a unique primitive word.
M\"obius inversion on this power decomposition gives
\(\sum_{k\mid\gcd(w,d)}\mu(k)\binom{w/k}{d/k}\)
primitive words, proving the formula. Finally, there are
\(\binom{w-1}{d-1}\) words ending in \(1\) in the component, which
gives the claimed product rank.
\end{proof}

\begin{table}[htbp]
\centering
\begin{tabular}{c|rrrrr}
\toprule
Weight \(w\)&\(d=1\)&\(d=2\)&\(d=3\)&\(d=4\)&\(d=5\)\\
\midrule
3&1&1&0&--&--\\
4&1&1&1&0&--\\
5&1&2&2&1&0\\
6&1&2&3&2&1\\
7&1&3&5&5&3\\
8&1&3&7&8&7\\
9&1&4&9&14&14\\
\bottomrule
\end{tabular}
\caption{The exact formal indecomposable counts \(\ell_{w,d}\).
They count Lyndon generators, not independent numerical constants.}
\label{tab:witt}
\end{table}

For \(w=6,d=3\), the ten words reduce modulo products to three Lyndon
words:
\[
000111,\qquad001011,\qquad001101,
\]
corresponding respectively to \(F_{4,1,1}\), \(F_{3,2,1}\),
and \(F_{3,1,2}\). This explains the manuscript's finite rank-seven
observation and gives its uniform extension.
The companion code records the full polynomial normal forms of all ten
triples. For a compact view of the quotient, put
\(X=[F_{4,1,1}],Y=[F_{3,2,1}],Z=[F_{3,1,2}]\), where brackets mean
classes modulo products. Exact triangular elimination gives
\[
\begin{array}{c|rrr@{\qquad}c|rrr}
 &X&Y&Z&&X&Y&Z\\ \hline
{}[F_{4,1,1}]&1&0&0 &[F_{1,4,1}]&6&3&1\\
{}[F_{3,2,1}]&0&1&0 &[F_{1,3,2}]&-12&-6&-2\\
{}[F_{3,1,2}]&0&0&1 &[F_{1,2,3}]&12&6&2\\
{}[F_{2,3,1}]&-9&-4&-1 &[F_{1,1,4}]&-6&-3&-1\\
{}[F_{2,2,2}]&12&4&0 &[F_{2,1,3}]&-3&-1&0
\end{array}
\]
For example, one full identity before taking the quotient is
\begin{equation}
F_{2,2,2}(z)=12F_{4,1,1}(z)+4F_{3,2,1}(z)
-2\Li_2(z)F_{3,1}(z)+\frac{\Li_2(z)^3}{6}.
\label{eq:full-222}
\end{equation}
Re-expanding its right side into words gives the single word \(010101\)
with coefficient one. This supplies a finite rational certificate,
including the product terms that a quotient table suppresses.

\subsection{An elementary all-weight double elimination}

At depth two the count simplifies to
\[
\ell_{w,2}=\left\lfloor\frac{w-1}{2}\right\rfloor.
\]
There is also a direct triangular proof that avoids general shuffle
algebra. For \(p+q=w\), shuffling the two single-polylogarithm words gives
\begin{align}
\Li_p(z)\Li_q(z)
={}&\sum_{j=1}^{p}\binom{w-j-1}{q-1}F_{w-j,j}(z)
+\sum_{j=1}^{q}\binom{w-j-1}{p-1}F_{w-j,j}(z).
\label{eq:general-double-shuffle}
\end{align}
The binomial coefficient counts the interleavings of the zero blocks
before the second nonzero letter.

\begin{proposition}[Independent double product rows]\label{prop:double-rank}
For fixed \(w\), the rows of
\eqref{eq:general-double-shuffle} with
\(1\le p\le\lfloor w/2\rfloor\) have rank
\(\lfloor w/2\rfloor\) on the \(w-1\) double coordinates.
They express the entire family through products of singles and the
\(\lfloor(w-1)/2\rfloor\) coordinates
\[
\{F_{w-j,j}:1\le j\le\lfloor(w-1)/2\rfloor\}.
\]
This set is empty when \(w=2\).
\end{proposition}

\begin{proof}
In the row indexed by \(p\), the largest inner index appearing is
\(q=w-p\). Its coefficient is \(1\) when \(p<q\), and \(2\)
when \(p=q\). No larger inner index occurs.
Order the rows with \(p\) decreasing and use the columns with
\(j=w-p\); the corresponding square block is triangular with nonzero
diagonal. All products of two depth-one words have already been included,
since swapping \(p,q\) repeats a row. This proves the rank and the
elimination statement.
\end{proof}

At \(z=i\), taking imaginary parts of these equations gives an
unconditional spanning bound for the \(g_{a,b}\), with explicit
single-value terms. It can be sharpened in the parity sector, as the
next section shows.

\subsection{The weight-five relation is two shuffles}

The words for \(\Li_1\Li_4\) and \(\Li_2\Li_3\) give
\begin{align}
\Li_1(z)\Li_4(z)
&=2F_{4,1}(z)+F_{3,2}(z)+F_{2,3}(z)+F_{1,4}(z),
\label{eq:shuffle14}\\
\Li_2(z)\Li_3(z)
&=6F_{4,1}(z)+3F_{3,2}(z)+F_{2,3}(z).
\label{eq:shuffle23}
\end{align}
The exact single values are
\[
\Li_1(i)=-\frac L2+\frac{\pi i}{4},\quad
\Li_2(i)=-\frac{\pi^2}{48}+iG,\quad
\Li_3(i)=-\frac{3\zeta(3)}{32}+\frac{\pi^3i}{32},
\]
\[
\Li_4(i)=-\frac{7\pi^4}{11520}+i\beta(4).
\]
Consequently
\begin{align}
2g_{41}+g_{32}+g_{23}+g_{14}
&=-\frac12\beta(4)L-\frac{7\pi^5}{46080},
\label{eq:wt5-first}\\
6g_{41}+3g_{32}+g_{23}
&=-\frac{\pi^5}{1536}-\frac{3G\zeta(3)}{32}.
\label{eq:wt5-second}
\end{align}
Here \(g_{41}\) means \(g_{4,1}\), and similarly for other subscripts.

\begin{corollary}[Proof and strengthening of the recorded relation]
\label{cor:wt5}
The manuscript's identity
\[
576g_{41}+288g_{32}+736g_{23}+960g_{14}
=21G\zeta(3)-480\beta(4)L
\]
holds. Moreover the four \(g\)-values are spanned by \(g_{41},g_{32}\)
and the single-value products \(\pi^5,G\zeta(3),\beta(4)L\).
Explicitly,
\begin{align}
g_{23}
&=-6g_{41}-3g_{32}-\frac{\pi^5}{1536}
  -\frac{3G\zeta(3)}{32},\label{eq:g23-elim}\\
g_{14}
&=4g_{41}+2g_{32}+\frac{23\pi^5}{46080}
  +\frac{3G\zeta(3)}{32}-\frac12\beta(4)L.
\label{eq:g14-elim}
\end{align}
\end{corollary}

\begin{proof}
Multiply \eqref{eq:wt5-first} by \(960\), multiply
\eqref{eq:wt5-second} by \(-224\), and add.
The \(\pi^5\) coefficients cancel exactly. Solving the two
independent equations instead gives \eqref{eq:g23-elim} and
\eqref{eq:g14-elim}.
\end{proof}

This proof explains why a search can return an apparently complicated
relation while missing the simpler structure: eliminating a chosen
single-value coordinate can enlarge coefficients.
The mathematical content is visible before the elimination.

\subsection{The surviving mixed relation}

For the harmonic sums used by the manuscript,
\[
S_p=\sum_{n\ge0}\frac{(-1)^nH_n}{(2n+1)^p},
\]
the exact bridge to colored doubles is
\begin{equation}
S_p=\Im\bigl(F_{p,1}(i)+\Li_{p,1}(i,-1)\bigr).
\label{eq:S-mixed}
\end{equation}
Indeed, for an odd outer index \(m=2n+1\), the inner coefficient
\(\sum_{k<m}(1+(-1)^k)/k\) is \(H_n\); even outer indices contribute
zero to the imaginary part. Absolute convergence justifies the argument
for \(p>1\), and Abel limits cover \(p=1\).

The recorded \(S_4\) candidate becomes, after the proved
\eqref{eq:g23-elim}, the more economical statement
\begin{conjecture}[A remaining mixed Gaussian reduction]\label{conj:S4}
\begin{equation}
S_4=\frac{58g_{41}+24g_{32}}7
       +\frac{19\pi^5}{3584}-2\beta(4)\log2.
\label{eq:S4-short}
\end{equation}
\end{conjecture}
This is equivalent to the existing candidate, not a newly discovered
independent relation. Independent integral evaluations support it,
but the present article does not supply its functional-equation proof.
In particular, it must not be silently added to the exact shuffle
certificate. It identifies a concrete additional colored relation to
seek. Neither its numerical support nor the dimensions of a selected
finite relation system establish arithmetic independence of any
surviving coordinate.

% Self-contained theorem/proof sections for integration into the parent report.
% Requires amsmath, amssymb, amsthm and a \Li macro.

\section{An explicit binomial parity reduction for every double weight}
\label{sec:double-parity}

Throughout this section the convention is
\[
 F_{a,b}(z)=\Li_{a,b}(z,1)
   =\sum_{n>m>0}\frac{z^n}{n^a m^b},\qquad
 U_n(z)=\operatorname{Re}\Li_n(z),\quad
 V_n(z)=\operatorname{Im}\Li_n(z).
\]
We take $|z|=1$, $z\ne1$, and use radial boundary values. These are also
the ordinary convergent series values: summation by parts applies to
$H_{n-1}^{(b)}/n^a$, whose total variation is finite and which tends to zero.
The result below is an elementary explicit specialization of the established
multiple-polylogarithm parity theorem. It should not be presented as a new
general parity theorem: Panzer proved the general theorem in 2017, including
explicit formulas at depths two and three; Umezawa supplied a general-depth
explicit formula in 2025 \cite{Panzer,Umezawa}. Its contribution here
is a short independently checkable formula in the manuscript's convention,
including the boundary cases with an index equal to one.

For $w\ge2$ define the parity component
\[
 E_{a,b}(z)=
 \begin{cases}
 \operatorname{Im}F_{a,b}(z),&w=a+b\text{ even},\\
 \operatorname{Re}F_{a,b}(z),&w=a+b\text{ odd},
 \end{cases}
\quad
 W_n(z)=
 \begin{cases}V_n(z),&w\text{ even},\\U_n(z),&w\text{ odd}.
 \end{cases}
\]
The choice of $W_n$ depends on the fixed total weight $w$, not on $n$.

\begin{theorem}[Explicit double parity formula]\label{thm:binomial-parity}
Let $w\ge2$, $1\le p\le w-1$, and $q=w-p$. Put
\begin{align*}
 A_p(z)&=\sum_{j=1}^{p}\binom{w-j-1}{q-1}E_{w-j,j}(z),\\
 P_p(z)&=
 \begin{cases}
 U_q(z)V_p(z),&w\text{ even},\ p\text{ odd},\\
 U_p(z)V_q(z),&w\text{ even},\ p\text{ even},\\
 -V_p(z)V_q(z),&w\text{ odd},\ p\text{ odd},\\
 U_p(z)U_q(z),&w\text{ odd},\ p\text{ even}.
 \end{cases}
\end{align*}
Then
\begin{align}
 A_p(z)={}&P_p(z)
 +\sum_{\substack{3\le j\le p\\j\text{ odd}}}
       \binom{w-j-1}{q-1}\zeta(j)W_{w-j}(z)
 -\sum_{\substack{3\le j\le q\\j\text{ odd}}}
       \binom{w-j-1}{p-1}\zeta(j)W_{w-j}(z)\notag\\
 &-\frac12\left[\binom{w-1}{q}-\binom{w-1}{p}\right]W_w(z)
 -\boldsymbol{1}_{w\text{ odd}}\frac{(-1)^p}{2}\zeta(w).
 \label{eq:parity-A}
\end{align}
Each individual double is recovered by the explicit inverse transform
\begin{equation}\label{eq:parity-inverse}
 E_{w-b,b}(z)=\sum_{p=1}^{b}(-1)^{b-p}
                  \binom{w-p-1}{b-p}A_p(z),\qquad 1\le b<w.
\end{equation}
Consequently every imaginary part of even total weight and every real part
of odd total weight in the same-argument double family is an explicit
polynomial in single polylogarithms and ordinary zeta values.
\end{theorem}

\begin{proof}
We supply the analytic identity responsible for the reduction, so that no
rank computation or numerical recognition enters the proof. Define
\[
 T_{p,q}(k)=\sum_{m=1}^{\infty}\frac1{(m+k)^p m^q},\qquad k\ge1.
\]
Partial fractions at $m=0$ and $m=-k$ give
\begin{align}
 \frac1{(m+k)^p m^q}
 ={}&\sum_{j=1}^{q}(-1)^{q-j}
       \binom{w-j-1}{p-1}\frac1{k^{w-j}m^j}\notag\\
 &+(-1)^q\sum_{j=1}^{p}
       \binom{w-j-1}{q-1}\frac1{k^{w-j}(m+k)^j}.
 \label{eq:parity-pf}
\end{align}
The two $j=1$ divergences cancel. Summing first to a finite cutoff and
then letting it tend to infinity yields
\begin{align}
 T_{p,q}(k)={}&
 \sum_{j=2}^{q}(-1)^{q-j}\binom{w-j-1}{p-1}
                \frac{\zeta(j)}{k^{w-j}}
 +(-1)^q\sum_{j=2}^{p}\binom{w-j-1}{q-1}
                \frac{\zeta(j)}{k^{w-j}}\notag\\
 &+(-1)^{q+1}\sum_{j=1}^{p}\binom{w-j-1}{q-1}
                \frac{H_k^{(j)}}{k^{w-j}}.
 \label{eq:parity-T}
\end{align}
Write
\[
 \mathcal A_p(z)=\sum_{j=1}^{p}\binom{w-j-1}{q-1}F_{w-j,j}(z),
 \qquad C_p=\binom{w-1}{q},
\]
and
\[
 \mathcal U_{p,q}(z)=
 \sum_{j=2}^{q}(-1)^{q-j}\binom{w-j-1}{p-1}
                 \zeta(j)\Li_{w-j}(z)
 +(-1)^q\sum_{j=2}^{p}\binom{w-j-1}{q-1}
                 \zeta(j)\Li_{w-j}(z).
\]
Using $H_k^{(j)}=H_{k-1}^{(j)}+k^{-j}$ and the hockey-stick identity,
equation~\eqref{eq:parity-T} becomes
\begin{equation}\label{eq:parity-R}
 R_{p,q}(z):=\sum_{k\ge1}z^kT_{p,q}(k)
 =\mathcal U_{p,q}(z)+(-1)^{q+1}
        \bigl(\mathcal A_p(z)+C_p\Li_w(z)\bigr).
\end{equation}

There are two product identities. The familiar same-argument shuffle is
\begin{equation}\label{eq:parity-shuffle}
 \Li_p(z)\Li_q(z)=\mathcal A_p(z)+\mathcal A_q(z).
\end{equation}
The Fourier-convolution identity is
\begin{equation}\label{eq:parity-convolution}
 \Li_p(z)\Li_q(z^{-1})
       =\zeta(w)+R_{p,q}(z)+R_{q,p}(z^{-1}).
\end{equation}
For completeness, this remains valid when $p=1$ or $q=1$. Each
$\Li_p(e^{i\theta})$ belongs to $L^2$ because $\sum n^{-2p}<\infty$.
The Fourier coefficients of their product are therefore the absolutely
convergent convolutions, equal to $\zeta(w)$ at index zero and to
$T_{p,q}(k)$ and $T_{q,p}(k)$ at positive and negative indices. Both
sequences $T_{p,q}(k)$ decrease to zero. Dirichlet's test gives convergence
of the two Fourier series, uniform on closed subarcs away from $z=1$.
Abel--Poisson summation of the $L^1$ product converges to the continuous
product at every such point, proving~\eqref{eq:parity-convolution} there.
Thus the derivation has never multiplied divergent zeta series.

If $w$ is even, take imaginary parts in
\eqref{eq:parity-shuffle}--\eqref{eq:parity-convolution}. Since $p$ and $q$
have the same parity, these give the sum and difference of $A_p,A_q$.
Moreover,
\begin{align*}
 &\operatorname{Im}\bigl(\mathcal U_{p,q}(z)+
                         \mathcal U_{q,p}(z^{-1})\bigr)\\
 &\quad=2(-1)^p\left[
  \sum_{\substack{3\le j\le p\\j\text{ odd}}}
       \binom{w-j-1}{q-1}\zeta(j)V_{w-j}(z)
 -\sum_{\substack{3\le j\le q\\j\text{ odd}}}
       \binom{w-j-1}{p-1}\zeta(j)V_{w-j}(z)\right].
\end{align*}
Solving the two elementary equations gives~\eqref{eq:parity-A} for even
$w$. If $w$ is odd, take real parts. Now $p,q$ have opposite parity and
the same two equations again give the sum and difference. The analogous
last display has $-2(-1)^p$ in place of $2(-1)^p$, with $U$ in place of
$V$. The constant term $\zeta(w)$ contributes
$-(-1)^p\zeta(w)/2$. This proves~\eqref{eq:parity-A} in the odd case too.

Finally the coefficient matrix in the definition of $A_p$ is lower
triangular with diagonal one. Its inverse follows from
\[
 \binom{w-j-1}{p-j}\binom{w-p-1}{b-p}
 =\binom{w-j-1}{b-j}\binom{b-j}{p-j}
\]
and $\sum_{p=j}^{b}(-1)^{b-p}\binom{b-j}{p-j}=\delta_{b,j}$.
This proves~\eqref{eq:parity-inverse}.
\end{proof}

\subsection{Consequences for Gaussian and Eisenstein values}
\label{subsec:gaussian-doubles}

At $z=i$ we have $V_n(i)=\beta(n)$,
$U_1(i)=-\log2/2$, and
\[
 U_n(i)=-2^{-n}(1-2^{1-n})\zeta(n),\qquad n\ge2.
\]
The odd beta values are rational multiples of powers of $\pi$.
Thus Theorem~\ref{thm:binomial-parity} supplies an explicit all-weight
reduction to the Gaussian single-value algebra in the appropriate parity.
It proves all five previously numerical weight-six equations:
\begin{align*}
2048\operatorname{Im}\Li_{5,1}(i,1)
 &=-64\pi^3\zeta(3)-527\pi\zeta(5)+4096\beta(6),\\
1536\operatorname{Im}\Li_{4,2}(i,1)
 &=96\pi^3\zeta(3)-32\pi^2\beta(4)+1581\pi\zeta(5)-8448\beta(6),\\
1024\operatorname{Im}\Li_{3,3}(i,1)
 &=-3\pi^3\zeta(3)+64\pi^2\beta(4)-1581\pi\zeta(5)+4608\beta(6),\\
23040\operatorname{Im}\Li_{2,4}(i,1)
 &=-14\pi^4\beta(2)+135\pi^3\zeta(3)-1440\pi^2\beta(4)\\
 &\hspace{17mm}+23715\pi\zeta(5)-69120\beta(6),\\
92160\operatorname{Im}\Li_{1,5}(i,1)
 &=-150\pi^5\log2+56\pi^4\beta(2)-270\pi^3\zeta(3)\\
 &\hspace{17mm}+1920\pi^2\beta(4)-675\pi\zeta(5).
\end{align*}
An especially simple all-order subfamily is
\begin{equation}\label{eq:leading-index-family}
 \operatorname{Im}\Li_{2m-1,1}(z,1)
  =(m-1)V_{2m}(z)+U_{2m-1}(z)V_1(z)
   -\sum_{k=1}^{m-1}\zeta(2k+1)V_{2m-2k-1}(z),
 \qquad m\ge1.
\end{equation}
For example, this gives the further proved weight-eight identity
\[
 \operatorname{Im}\Li_{7,1}(i,1)
 =3\beta(8)-\frac{5\pi^5}{1536}\zeta(3)
             -\frac{\pi^3}{32}\zeta(5)
             -\frac{8255\pi}{32768}\zeta(7).
\]
At $z=\rho=e^{2\pi i/3}$ write $C_n=\operatorname{Im}\Li_n(\rho)$.
The same theorem applies unchanged; in particular
\[
 \operatorname{Im}\Li_{5,1}(\rho,1)
 =2C_6-\frac{2\pi^3}{81}\zeta(3)-\frac{121\pi}{486}\zeta(5),
\qquad
 \operatorname{Im}\Li_{3,1}(\rho,1)
 =C_4-\frac{13\pi}{54}\zeta(3).
\]
These statements assert reduction identities, not numerical or motivic
independence of the remaining single values.

\section{Reconciliation of the inverse-argument mixed values}
\label{sec:mixed-reconciliation}

The following results rederive conclusions already established in the
project's 8 October 2026 continuation. They are included because the
current consolidated manuscript still reports the corresponding values
as unexplained candidate directions. They must not be counted as newly
resolved conjectures in this continuation.

For $|z|=1$, $z\ne1$, the same function introduced in the Fourier proof is
precisely the inverse-argument mixed value:
\begin{equation}\label{eq:mixed-R-identification}
 R_{a,b}(z)=\Li_{a,b}(z,z^{-1}).
\end{equation}
Indeed, put $n=m+k$ in the defining series. To justify the boundary
rearrangement when an index is one, first insert a radial factor $r^n$ in
the original outer sum. The resulting positive coefficient at difference
$k$ is
$r^k\sum_{m\ge1}r^m/((m+k)^a m^b)$; it decreases in $k$ and is bounded
by $T_{a,b}(k)$. Uniform Dirichlet tail bounds and the finite initial
segment then permit $r\uparrow1$, giving~\eqref{eq:mixed-R-identification}.

Let $\mathcal M_{a,b}$ denote the imaginary part of the mixed value at even
weight and its real part at odd weight. With the notation of
Theorem~\ref{thm:binomial-parity}, the explicit reduction is
\begin{align}\label{eq:mixed-parity-explicit}
 \mathcal M_{a,b}(z)
 ={}&(-1)^{b+1}\left[
 P_a(z)+\frac12\binom{w}{a}W_w(z)
 -\sum_{\substack{2\le j\le a\\j\text{ even}}}
     \binom{w-j-1}{b-1}\zeta(j)W_{w-j}(z)\right.\notag\\
 &\left.\hspace{24mm}
 -\sum_{\substack{2\le j\le b\\j\text{ even}}}
     \binom{w-j-1}{a-1}\zeta(j)W_{w-j}(z)\right]
 -\boldsymbol{1}_{w\text{ odd}}\frac{\zeta(w)}2,
 \qquad w=a+b.
\end{align}
This follows directly by substituting~\eqref{eq:parity-A} into
\eqref{eq:parity-R}: the odd-zeta terms cancel and
$\binom{w-1}{a}+\binom{w-1}{b}=\binom{w}{a}$.
In particular,
\begin{align*}
 \operatorname{Re}\Li_{2m,1}(z,z^{-1})
 &=\frac{2m+1}{2}U_{2m+1}(z)-\frac12\zeta(2m+1)
    +U_1(z)U_{2m}(z)\\
 &\quad-\sum_{k=1}^{m}\zeta(2k)U_{2m+1-2k}(z),\\
 \operatorname{Im}\Li_{2m-1,1}(z,z^{-1})
 &=mV_{2m}(z)+U_1(z)V_{2m-1}(z)
    -\sum_{k=1}^{m-1}\zeta(2k)V_{2m-2k}(z).
\end{align*}
At the manuscript's four proposed mixed directions these give
\begin{align*}
 \operatorname{Re}\Li_{4,1}(i,-i)
 &=\frac{3\pi^4\log2}{512}+\frac{\pi^2\zeta(3)}{64}
                               -\frac{587\zeta(5)}{1024},\\
 \operatorname{Im}\Li_{5,1}(i,-i)
 &=3\beta(6)-\frac{5\pi^5\log2}{3072}
             -\frac{\pi^2\beta(4)}6-\frac{\pi^4G}{90},\\
 \operatorname{Re}\Li_{4,1}(\rho^2,\rho)
 &=\frac{2\pi^4\log3}{243}+\frac{2\pi^2\zeta(3)}{27}
                               -\frac{281\zeta(5)}{162},\\
 \operatorname{Im}\Li_{5,1}(\rho^2,\rho)
 &=-3C_6+\frac{\pi^5\log3}{729}
             +\frac{\pi^2C_4}{6}+\frac{\pi^4C_2}{90}.
\end{align*}
Both weight-five real reductions have integral relation height below
$10^4$: clearing denominators gives coefficients
$(1024,-6,-16,587)$ and $(486,-4,-36,843)$, respectively.
Consequently their reported non-reduction against the stated complete
single-value basket is not evidence of a large-coefficient obstruction;
the value generation, basis construction, or search setup requires an
audit. The rigorous formulas settle the mathematical status regardless
of that computational history.

\section{Classical reductions for every index with one two}
\label{sec:one-two}

This section uses the one-variable abbreviation
$\Li_{k_1,\ldots,k_d}(z)=\Li_{k_1,\ldots,k_d}(z,1,\ldots,1)$.
Let $\{1\}^{r}$ denote $r$ consecutive ones. Work first on $0<z<1$, and
then continue analytically along paths avoiding the cuts of the displayed
functions. In particular, the formulas hold at $z=i$ by continuation
through the upper half-plane, with the principal values used below.

\begin{theorem}[One two at arbitrary depth]\label{thm:one-two}
For $d\ge1$, putting $a=\log(1-z)$ gives
\begin{equation}\label{eq:height-one-general}
 \Li_{2,\{1\}^{d-1}}(z)
 =\zeta(d+1)+\frac{(-1)^d a^d\log z}{d!}
   +\sum_{k=0}^{d-1}\frac{(-1)^{k+1}a^k}{k!}
                       \Li_{d+1-k}(1-z).
\end{equation}
For all $r,s\ge0$, put $\alpha=-\log(1-z)$. Then
\begin{equation}\label{eq:one-two-any-position}
 \Li_{\{1\}^{r},2,\{1\}^{s}}(z)
 =\sum_{j=0}^{r}(-1)^j\binom{s+j+1}{j}
      \frac{\alpha^{r-j}}{(r-j)!}
                     \Li_{2,\{1\}^{s+j}}(z).
\end{equation}
Consequently every value with exactly one index equal to two and all other
indices equal to one reduces to classical polylogarithms and logarithms.
\end{theorem}

The height-one functional equation is classical; a modern treatment and
arithmetic analogue appear in \cite{NakamuraShiraishi}. The proof below
also records the shuffle reduction when the two occupies any position.

\begin{proof}
The iterated integral for $d$ ones is
$\Li_{\{1\}^{d}}(z)=\alpha^d/d!$. Hence
\[
 \Li_{2,\{1\}^{d-1}}(z)
 =\frac{(-1)^d}{d!}\int_0^z\frac{\log^d(1-t)}{t}\,dt.
\]
Set $s=1-t$ and integrate by parts $d$ times. The antiderivative
\[
 -\log^d s\log(1-s)
  +\sum_{j=1}^{d}(-1)^j\frac{d!}{(d-j)!}
                       \log^{d-j}s\Li_{j+1}(s)
\]
has derivative $\log^d(s)/(1-s)$. At $s=1$ its limiting value is
$(-1)^d d!\zeta(d+1)$. This gives~\eqref{eq:height-one-general} including
the constant term and its sign.

For the second identity, locate the unique zero letter in the iterated
integral word $1^r0\,1^{s+1}$. Its variable is $t$. Integrating the
$r$ outer and $s+1$ inner identical letters gives
\[
 \Li_{\{1\}^{r},2,\{1\}^{s}}(z)
 =\frac1{r!(s+1)!}\int_0^z
    \bigl(\alpha(z)-\alpha(t)\bigr)^r\alpha(t)^{s+1}\frac{dt}{t}.
\]
Expanding the first factor and using the preceding integral formula
proves~\eqref{eq:one-two-any-position}. Analytic continuation then gives
the asserted complex values.
\end{proof}

\subsection{Proof of all three weight-four Gaussian triple formulas}
\label{subsec:gaussian-triples}

Put $L=\log2$, $G=\beta(2)$, $r=(1+i)/2$, and
$\lambda_n=\operatorname{Im}\Li_n(r)$,
$\mu_n=\operatorname{Re}\Li_n(r)$. The elementary dilogarithm reflection,
Landen transformation, and trilogarithm reflection give
\begin{align}
 \mu_2&=\frac{5\pi^2}{96}-\frac{L^2}{8},&
 \lambda_2&=G-\frac{\pi L}{8},\notag\\
 \mu_3&=\frac{35}{64}\zeta(3)-\frac{5\pi^2L}{192}
                       +\frac{L^3}{48}.
 \label{eq:halfpoint-auxiliary}
\end{align}
For clarity, the trilogarithm identity used for the last formula is
\begin{align*}
 &\Li_3(z)+\Li_3(1-z)+\Li_3\!\left(-\frac{z}{1-z}\right)\\
 &\quad=\zeta(3)+\zeta(2)\log(1-z)
  -\frac12\log z\log^2(1-z)+\frac16\log^3(1-z).
\end{align*}
It follows by differentiation using the dilogarithm identities and fixing
the constant as $z\to0$; it is then continued to $z=r$, where
$1-r=\bar r$ and $-r/(1-r)=-i$. Taking real parts proves the displayed
$\mu_3$ evaluation. The first two formulas follow respectively from
$\Li_2(r)+\Li_2(\bar r)=\zeta(2)-\log r\log\bar r$ and
$\Li_2(i)+\Li_2(\bar r)=-\tfrac12\log^2(1-i)$.

To evaluate the height-one formulas at $i$, write
\[
 a=\log(1-i)=\frac L2-\frac{\pi i}{4},\qquad
 c=\log(-(1-i))=\frac L2+\frac{3\pi i}{4}.
\]
Since $(1-i)^{-1}=r$, the needed principal-branch inversion formulas are
\begin{align*}
 \Li_2(1-i)&=-\Li_2(r)-\frac{\pi^2}{6}-\frac{c^2}{2},\\
 \Li_3(1-i)&=\Li_3(r)-\frac{\pi^2c}{6}-\frac{c^3}{6},\\
 \Li_4(1-i)&=-\Li_4(r)-\frac{7\pi^4}{360}
                         -\frac{\pi^2c^2}{12}-\frac{c^4}{24}.
\end{align*}
Substitution into~\eqref{eq:height-one-general} with $d=2,3$, using
$\log i=\pi i/2$, yields
\begin{align*}
 \operatorname{Re}\Li_{2,1}(i)&=\frac{29}{64}\zeta(3)-\frac{\pi G}{4},\\
 \operatorname{Im}\Li_{2,1}(i)&=-\frac{GL}{2}
                   +\frac{\pi L^2}{32}-\lambda_3+\frac{3\pi^3}{128},\\
 \operatorname{Im}\Li_{2,1,1}(i)&=
 \lambda_4+\frac{L\lambda_3}{2}
 -\frac{G(\pi^2-4L^2)}{32}
 -\frac{\pi(8L^3+105\zeta(3))}{768}.
\end{align*}
The real-part evaluation of $\Li_{2,1}(i)$ is also printed in
Panzer \cite{Panzer}, after reversing his index convention.
The remaining two identities follow from exact shuffle relations:
\begin{align*}
 \Li_1(z)\Li_{2,1}(z)
    &=\Li_{1,2,1}(z)+3\Li_{2,1,1}(z),\\
 \Li_2(z)\Li_{1,1}(z)
    &=3\Li_{2,1,1}(z)+2\Li_{1,2,1}(z)+\Li_{1,1,2}(z).
\end{align*}
Using $\Li_1(i)=-a$, $\Li_{1,1}(i)=a^2/2$, and
$\Li_2(i)=-\pi^2/48+iG$, one obtains
\begin{align*}
 \operatorname{Im}\Li_{1,2,1}(i)
 &=-3\lambda_4-L\lambda_3
   +\frac{G(\pi^2-4L^2)}{32}
   -\frac{3\pi^3L}{256}
   +\frac{\pi(2L^3+67\zeta(3))}{128},\\
 \operatorname{Im}\Li_{1,1,2}(i)
 &=3\lambda_4+\frac{L\lambda_3}{2}
   +\frac{5\pi^3L}{192}-\frac{163\pi\zeta(3)}{256}.
\end{align*}
These are exactly the three numerical formulas in the current manuscript,
now proved. The arbitrary-depth theorem proves additional such identities
without a new integer-relation search. It does not address the independent
status of the weight-six triples $(3,1,2),(3,2,1),(4,1,1)$, which contain
more than one zero letter and lie outside this reduction mechanism.


% Self-contained contribution for the Polylogarithms continuation article.
% Required packages: amsmath, amssymb, amsthm, booktabs, hyperref.
% Required theorem environments: theorem, lemma, corollary, remark.

\section{Proof certificates for the cubic and quartic trilogarithm ladders}
\label{sec:proved-small-ladders}

Chapter~3 of the audited manuscript records two trilogarithm combinations
at the positive roots of $r^3+r=1$ and $r^4+r=1$, with substantial numerical
evidence but without a local derivation from functional equations.
We give such derivations. The quartic identity belongs to a uniform
family; the cubic identity has a short certificate using two
specializations of Kummer's nine-term relation and four Rogers pentagons.
These are exact consequences of classical functional equations, not
claims that the underlying equations are new.

\subsection{A uniform ladder from a complementary pair of powers}

For $0<x<1$ the real Landen relation for the trilogarithm is
\begin{align}
\operatorname{Li}_3(x)+\operatorname{Li}_3(1-x)
 +\operatorname{Li}_3(1-x^{-1})
 &=\zeta(3)+\frac{\pi^2}{6}\log x
 -\frac12\log^2x\log(1-x)+\frac16\log^3x.
\label{eq:real-trilog-landen}
\end{align}
For a modern account and its arithmetic analogues see
Nakamura and Shiraishi, \emph{Landen's trilogarithm functional equation
and $\ell$-adic Galois multiple polylogarithms} \cite{NakamuraShiraishi}.
For completeness, a real-variable proof is especially short.
Differentiation of the left side, followed by
\[
\operatorname{Li}_2(x)+\operatorname{Li}_2(1-x)
 =\frac{\pi^2}{6}-\log x\log(1-x),\qquad
\operatorname{Li}_2(1-x^{-1})
 =-\operatorname{Li}_2(1-x)-\frac12\log^2x,
\]
gives
\[
\frac{\pi^2/6-\log x\log(1-x)}{x}
  +\frac{\log^2x}{2x(1-x)}.
\]
This is the derivative of the right side. The difference tends to zero
as $x\to1^-$, proving \eqref{eq:real-trilog-landen}. The two dilogarithm
relations used here follow themselves by one differentiation and an
endpoint value. All logarithms in the calculation are real, so there
is no branch ambiguity.

\begin{theorem}[Complementary-power trilogarithm ladder]
\label{thm:complementary-powers}
Let $a,b$ be positive integers with $a\le b$, and let $0<\rho<1$ be
the unique positive solution of $\rho^a+\rho^b=1$. Set
$\lambda=\log\rho$ and $d=b-a$. Then
\begin{align}
&\operatorname{Li}_3(\rho^a)+\operatorname{Li}_3(\rho^b)
 -\operatorname{Li}_3(\rho^d)
 +\frac14\operatorname{Li}_3(\rho^{2d})\notag\\
&\hspace{20mm}=\zeta(3)+\frac{a\pi^2}{6}\lambda
 +\frac{a^2(a-3b)}{6}\lambda^3.
\label{eq:complementary-powers}
\end{align}
When $a=b$, the powers $\rho^0$ are interpreted as $1$.
\end{theorem}

\begin{proof}
The polynomial $t^a+t^b$ is strictly increasing from $0$ to $2$ on
$[0,1]$, which proves existence and uniqueness. Substitute $x=\rho^a$
in \eqref{eq:real-trilog-landen}; its other arguments are
$1-x=\rho^b$ and $1-x^{-1}=-\rho^{b-a}$. The absolutely convergent
series identity
\[
\operatorname{Li}_3(z)+\operatorname{Li}_3(-z)
 =\frac14\operatorname{Li}_3(z^2),\qquad 0\le z\le1,
\]
eliminates the negative argument. The logarithmic terms become exactly
the right side of \eqref{eq:complementary-powers}.
\end{proof}

In particular, if $r_m^m+r_m=1$ and $\lambda_m=\log r_m$, then
\begin{align}
\operatorname{Li}_3(r_m)+\operatorname{Li}_3(r_m^m)
 -\operatorname{Li}_3(r_m^{m-1})
 +\frac14\operatorname{Li}_3(r_m^{2m-2})
 =\zeta(3)+\frac{\pi^2}{6}\lambda_m
 -\frac{3m-1}{6}\lambda_m^3.
\label{eq:all-m-trilog}
\end{align}
This supplies an infinite family of finite exact reductions, although
different pairs $(a,b)$ can describe the same base after a power
substitution. It does not assert that these exhaust all ladders for
any one number field.

\begin{corollary}[The quartic identity]
If $r^4+r=1$ and $\lambda=\log r$, then
\begin{align}
-12\operatorname{Li}_3(r)+12\operatorname{Li}_3(r^3)
 -12\operatorname{Li}_3(r^4)-3\operatorname{Li}_3(r^6)
 =22\lambda^3-2\pi^2\lambda-12\zeta(3).
\label{eq:quartic-proved}
\end{align}
\end{corollary}
\begin{proof}
Take $(a,b)=(1,4)$ in Theorem~\ref{thm:complementary-powers}
and multiply by $-12$.
\end{proof}

For $m=2$, \eqref{eq:all-m-trilog} reduces to the familiar golden value
\[
\operatorname{Li}_3(r_2^2)
 =\frac45\zeta(3)+\frac{2\pi^2}{15}\log r_2
 -\frac23\log^3r_2.
\]
For $m=3$ it gives a useful additional relation on the arguments
$r,r^2,r^3,r^4$; the manuscript's cubic identity instead uses
$r,r^2,r^3,r^6$. The change of the last argument requires more than
the one-variable Landen relation.

\subsection{An elementary dilogarithm certificate in the cubic field}

Throughout this subsection let $r^3+r=1$, $0<r<1$, and put
$\lambda=\log r$, $A=\log(1+r)$. Direct algebra in
$\mathbb Q[r]/(r^3+r-1)$ gives
\begin{align}
1-r&=r^3,&1-r^2&=r^3(1+r),&1-r^3&=r,\\
1-r^4&=r^2(1+r),&1+r^3&=r^2(1+r)^2,
 &1-r^6&=r^3(1+r)^2.
\label{eq:cubic-factorizations}
\end{align}
Define the Rogers dilogarithm on $(0,1)$ by
\[
R(x)=\operatorname{Li}_2(x)+\frac12\log x\log(1-x).
\]
It satisfies $R(x)+R(1-x)=\zeta(2)$ and the pentagon
\begin{align}
R(x)+R(y)
 =R(xy)+R\!\left(\frac{x(1-y)}{1-xy}\right)
  +R\!\left(\frac{y(1-x)}{1-xy}\right),\qquad 0<x,y<1.
\label{eq:rogers-pentagon}
\end{align}
The latter is a real-domain form of the classical five-term relation.
It may also be verified by differentiating in $x$ with fixed $y$ and
using the limit at $x=0$.

\begin{lemma}[Cubic dilogarithm identity]
\label{lem:cubic-dilog}
\begin{equation}
\operatorname{Li}_2(r^6)
 =6\operatorname{Li}_2(r^2)-2\operatorname{Li}_2(r)
  -4\operatorname{Li}_2(r^3).
\label{eq:cubic-dilog}
\end{equation}
\end{lemma}

\begin{proof}
Set $s=r/(1+r)$ and $t=r^4/(1+r)$, so every argument below lies
in $(0,1)$. Write $R_k=R(r^k)$ and $R_s=R(s)$, $R_t=R(t)$.
Applying \eqref{eq:rogers-pentagon} at $(x,y)=(r,r)$ gives
\[
2R_1=R_2+2R_s.
\]
At $(x,y)=(r,r^3)$ its two remaining arguments are $1-s$ and $t$.
Since $R_1+R_3=\zeta(2)$, reflection gives
\[
R_t=R_s-R_4.
\]
At $(x,y)=(r^2,r^4)$ its two remaining arguments are $s$ and $t$,
and therefore
\[
R_6=R_2+2R_4-2R_s.
\]
Finally, at $(x,y)=(r,r^2)$ these arguments are $1-r^2$ and $r^4$.
Using reflection twice gives
\[
R_4=2R_2-2R_3.
\]
Eliminating $R_s,R_t,R_4$ now yields
$R_6=6R_2-2R_1-4R_3$. The difference between $R(r^k)$ and
$\operatorname{Li}_2(r^k)$ is
$(k\lambda/2)\log(1-r^k)$. In the indicated combination these
terms cancel: by \eqref{eq:cubic-factorizations}, the required check is
\[
3\lambda(3\lambda+2A)
 -6\lambda(3\lambda+A)+3\lambda^2+6\lambda^2=0.
\]
This proves the claimed ordinary dilogarithm identity.
\end{proof}

\subsection{Two Kummer substitutions prove the cubic trilogarithm}

The single-valued trilogarithm used in the manuscript is
\[
P_3(z)=\operatorname{Re}\operatorname{Li}_3(z)
 -\log|z|\operatorname{Re}\operatorname{Li}_2(z)
 -\frac13\log^2|z|\log|1-z|.
\]
The value at $z=1$ is defined by continuity. It obeys
\[
P_3(z^{-1})=P_3(z),\qquad
P_3(z)+P_3(-z)=\frac14P_3(z^2),
\]
and the single-valued Landen relation
$P_3(x)+P_3(1-x)+P_3(1-x^{-1})=\zeta(3)$ for real $0<x<1$.
The latter follows from \eqref{eq:real-trilog-landen} and the two
dilogarithm identities in its proof.

We use Kummer's classical equation in the form stated and justified
by Zagier, \emph{Special values and functional equations of
polylogarithms}, \S7, Example~3, equation~(7) \cite{Zagier}:
\begin{align}
&2P_3(x)+2P_3(y)
 +2P_3\!\left(\frac{x(1-y)}{x-1}\right)
 +2P_3\!\left(\frac{y(1-x)}{y-1}\right)\notag\\
&\quad+2P_3\!\left(\frac{1-x}{1-y}\right)
 +2P_3\!\left(\frac{x(1-y)}{y(1-x)}\right)
 -P_3(xy)-P_3(x/y)\notag\\
&\quad-P_3\!\left(\frac{x(1-y)^2}{y(1-x)^2}\right)=2\zeta(3).
\label{eq:kummer-P3}
\end{align}
We only use specializations with nonzero denominators and finite
arguments. Since $P_3$ is single-valued, some real arguments exceeding
$1$ cause no branch correction in this equation; after inversion and
duplication, all remaining arguments are $r^k\in(0,1)$.

\begin{theorem}[The cubic identity]
\label{thm:cubic-proved}
Let $r^3+r=1$ with $0<r<1$, and let $\lambda=\log r$. Then
\begin{align}
&-60\operatorname{Li}_3(r)+54\operatorname{Li}_3(r^2)
 -32\operatorname{Li}_3(r^3)-3\operatorname{Li}_3(r^6)\notag\\
&\hspace{20mm}=24\lambda^3-4\pi^2\lambda-36\zeta(3).
\label{eq:cubic-proved}
\end{align}
\end{theorem}

\begin{proof}
Write $P_k=P_3(r^k)$. Landen at $r$ and duplication give
\[
\mathsf A:=P_1-P_2+P_3+\tfrac14P_4=\zeta(3).
\]
The nine arguments of \eqref{eq:kummer-P3} at
$(x,y)=(r,r^3)$ and $(-r,r^2)$, in their displayed order, are
\[
\begin{array}{c|rrrrrrrrr}
 &1&2&3&4&5&6&7&8&9\\ \hline
(r,r^3)&r&r^3&-r^{-1}&-r^5&r^2&r^{-4}&r^4&r^{-2}&r^{-6}\\
(-r,r^2)&-r&r^2&r^4&-r^{-1}&r^{-3}&-r^2&-r^3&-r^{-1}&-r^5
\end{array}
\]
respectively. Using inversion and
$P_3(-r^k)=\tfrac14P_{2|k|}-P_{|k|}$, these two equations reduce to
\begin{align*}
\mathsf B&:=\tfrac32P_2+2P_3+P_4-2P_5-P_6+\tfrac12P_{10}
 =2\zeta(3),\\
\mathsf C&:=-3P_1+\tfrac34P_2+3P_3+\tfrac52P_4+P_5
 -\tfrac14P_6-\tfrac14P_{10}=2\zeta(3).
\end{align*}
The exact rational row operation $4\mathsf C+2\mathsf B-48\mathsf A$
therefore proves
\begin{equation}
-60P_1+54P_2-32P_3-3P_6=-36\zeta(3).
\label{eq:cubic-P3-certificate}
\end{equation}

To recover ordinary trilogarithms, put
$c_1=-60,c_2=54,c_3=-32,c_6=-3$ and
\[
D=\sum_{k\in\{1,2,3,6\}}c_k\operatorname{Li}_3(r^k),\quad
S=\sum_{k\in\{1,2,3,6\}}c_k k\operatorname{Li}_2(r^k),\quad
T=\sum_{k\in\{1,2,3,6\}}c_k k^2\log(1-r^k).
\]
Expansion of the definition gives
$\sum c_kP_k=D-\lambda S-\lambda^2 T/3$.
Lemma~\ref{lem:cubic-dilog}, together with
$\operatorname{Li}_2(r)+\operatorname{Li}_2(r^3)
=\pi^2/6-3\lambda^2$, gives
\[
S=72\lambda^2-4\pi^2.
\]
The factorizations \eqref{eq:cubic-factorizations} give
\[
T=-60(3\lambda)+216(3\lambda+A)
 -288\lambda-108(3\lambda+2A)=-144\lambda.
\]
Consequently \eqref{eq:cubic-P3-certificate} becomes
$D-24\lambda^3+4\pi^2\lambda=-36\zeta(3)$, which is
\eqref{eq:cubic-proved}.
\end{proof}

\subsection{What the certificate establishes and what remains open}

The proofs settle the two displayed cubic and quartic trilogarithm
identities without period-independence assumptions, numerical recognition,
or unproved higher Bloch-group assertions. The universal family
\eqref{eq:complementary-powers} also gives a principled generator of
additional exact identities. It does not justify pseudointegration to
weight four or above: differentiating with respect to an algebraic base
does not preserve its isolated defining relation.

Natural next problems are to derive comparable local certificates for the
weight-five through weight-nine golden formulas, find the smallest useful
functional-equation closure for a chosen base, and extend the exact
argument/substitution checker to arbitrary algebraic number fields.
Any claim of minimality of the surviving period basket remains distinct
from the existence of the displayed reductions.


% Analytic research contribution.  The parent article supplies theorem environments.
% Compiled mathematics and provenance: manuscript commit afed07429d3d37eceb6c8e9e54cf4da2d3f39d53.
\section{Log-gamma moments: poles, late coefficients, and effective remainders}
\label{sec:gamma}
Let
\[
 M_n=\int_0^1\log^n\Gamma(x)\,dx,\qquad L=\log(2\pi),\qquad A=\gamma+L.
\]
The cubic moment in Chapter~07 is an appropriate test of whether a numerical
search has selected the right class of constants.  Its representation by
Tornheim derivatives is, however, already known.  Bailey, D.~Borwein and
J.~M.~Borwein explicitly evaluated $M_3$ in their Theorem~6 and $M_4$ in
Theorems~7--8 \cite{BBB}.  The manuscript's statement that discovering a
named atom for $M_3$ remains open should therefore be replaced by the
narrower question of reduction of those derivative coordinates in a specified
smaller algebra.  No independence result follows from the unsuccessful
single-zeta search.

Similarly, the full fixed-order expansion of $M_n/n!$ into the exponential
scales $k^{-n}$ is Theorem~8.1 of Amdeberhan et al.\ \cite{ACEMKM}.
The results proved below add an explicit late-coefficient asymptotic, identify
the divergent character of that expansion, and give a uniform remainder
bound when the number of retained terms grows linearly with $n$.
These are refinements of a known expansion, not a claim to have discovered
that expansion or the cubic Tornheim representation.

\subsection{A compact form of the cubic identity}
Define the holomorphic germ near $(1,1,1)$
\[
 T(r,s,t)=\sum_{m,n\ge1}\frac{1}{m^r n^s(m+n)^t}.
\]
A subscript $j$ denotes partial differentiation in its $j$th argument;
all unmarked Tornheim terms in the next formula are evaluated at $(1,1,1)$.
The double sums and all the derivatives used here converge absolutely.

\begin{theorem}[Cubic moment in one Tornheim differential operator]
\label{thm:cubic-compact}
With $A=\gamma+L$,
\begin{align}
 M_3={}&\frac{L^3}{8}+\frac{\pi^2L}{32}+\frac{3\zeta(3)}{16}
 +\frac{3L}{4\pi^2}
 \bigl(A^2\zeta(2)-2A\zeta'(2)+\zeta''(2)\bigr)\notag\\
 &+\frac{3}{8\pi^2}
 \left(A^2T-2AT_3+2T_{13}-T_{12}\right).
 \label{eq:cubic-compact}
\end{align}
This is an equivalent form of the known cubic evaluation in
\cite[Theorem~6]{BBB}.
\end{theorem}

\begin{proof}
Write the Kummer Fourier series as
\[
 \log\Gamma(x)=\frac L2+\frac12 C(x)+\frac1\pi S(x),
 \quad C(x)=\sum_{j\ge1}\frac{\cos(2\pi jx)}j,
 \quad S(x)=\sum_{j\ge1}\frac{A+\log j}{j}\sin(2\pi jx).
\]
Reflection $x\mapsto1-x$ makes $C$ even and $S$ odd.  Parseval gives
\[
 \int_0^1 C^2=\frac12\zeta(2),\qquad
 \int_0^1 S^2=\frac12\bigl(A^2\zeta(2)-2A\zeta'(2)+\zeta''(2)\bigr).
\]
The cosine resonance gives $\int C^3=3\zeta(3)/2$.  Equivalently, one can
differentiate the beta integral for $\int_0^1(2\sin\pi x)^u dx$ three times
at zero and use $C=-\log(2\sin\pi x)$.
For the only remaining cubic term, the elementary product rule for two sines
and one cosine gives
\begin{align*}
 \int_0^1 CS^2
 &=\frac12\sum_{m,n\ge1}
   \frac{(A+\log m)(A+\log(m+n))}{mn(m+n)}\\
 &\hspace{1em}-\frac14\sum_{m,n\ge1}
   \frac{(A+\log m)(A+\log n)}{mn(m+n)}\\
 &=\frac14\left(A^2T-2AT_3+2T_{13}-T_{12}\right).
\end{align*}
Expansion of the cube now proves the formula.

For the convergence justification, first put an Abel factor $r^j$ in every
Fourier series, $0<r<1$.  All products can then be integrated termwise.
The functions $C$ and $S$ belong to every finite $L^p(0,1)$, since the
identities with $\log\Gamma(x)$ and $\log\Gamma(1-x)$ bound them by
endpoint logarithms.  Poisson convolution therefore converges to them in
$L^3$, so the integrated cubic products converge.  In the resulting
resonance sums the logarithmic numerator is bounded by a constant times
$(1+\log(m+n))^2$.  Summing first over $m+n=q$ uses
\[
 \sum_{m=1}^{q-1}\frac{1}{m(q-m)q}=\frac{2H_{q-1}}{q^2},
\]
which proves absolute convergence and legitimizes the Abel limit.
\end{proof}

The double-coordinate derivative in \eqref{eq:cubic-compact} can also be
written as a positive sum.  Put
\[
 \mathcal R_A=\sum_{m,n\ge1}
 \frac{(A+\log(m+n))^2-
 \log((m+n)/m)\log((m+n)/n)}{mn(m+n)}.
\]
Symmetry in $m,n$ proves
$\mathcal R_A=A^2T-2AT_3+2T_{13}-T_{12}$.
The numerator is positive: after expansion it is
$A^2+2A\log(m+n)+(\log m+\log n)\log(m+n)-\log m\log n$.
It is bounded above by $(A+\log(m+n))^2$.
Consequently, if $\mathcal R_{A,N}$ restricts $m+n\le N$ and $N\ge3$ is an integer, then
\begin{equation}
 0<\mathcal R_A-\mathcal R_{A,N}
 \le\frac2N\bigl(P(\log N)+P'(\log N)+P''(\log N)+P'''(\log N)\bigr),
 \quad P(u)=(1+u)(A+u)^2.
 \label{eq:tornheim-tail}
\end{equation}
Indeed $H_{q-1}\le1+\log q$, and
$(1+\log x)(A+\log x)^2/x^2$ decreases for $x\ge3$.
Integration of the resulting tail gives \eqref{eq:tornheim-tail}.
This bound is elementary and rigorous but converges only as
$O(\log^3N/N)$; it is useful as an independent sign and normalization check,
not as the preferred high-precision algorithm.

\subsection{Meromorphic moment generating function}
For $\Re z<1$, let
\[
 F(z)=\int_0^1\Gamma(x)^z\,dx
     =\sum_{n\ge0}\frac{M_n}{n!}z^n\quad (|z|<1).
\]
The power uses the real logarithm of $\Gamma(x)>0$.  Introduce the positive
numbers
\begin{equation}
 c_k=\frac1k[u^{k-1}]\Gamma(1-u)^k,\qquad k\ge1.
 \label{eq:ck-lagrange}
\end{equation}
Their first values are
\[
 c_1=1,\quad c_2=\gamma,\quad
 c_3=\frac{3\gamma^2+\zeta(2)}2,\quad
 c_4=\frac{8\gamma^3+6\gamma\zeta(2)+\zeta(3)}3.
\]
These agree with the coefficients in \cite[Theorem~8.1]{ACEMKM}.

\begin{proposition}[All moment poles and their residues]
\label{prop:moment-poles}
The function $F$ extends meromorphically to the whole $z$-plane.  Its only
poles are simple poles at every positive integer, with
\[
 \operatorname*{Res}_{z=k}F(z)=(-1)^k k c_k\ne0.
\]
\end{proposition}
\begin{proof}
Fix $0<a<1$ and expand
\[
 \Gamma(x)^z=x^{-z}\Gamma(1+x)^z
 =x^{-z}\sum_{j=0}^{J-1}b_j(z)x^j+x^{J-z}E_J(x,z).
\]
The $b_j$ are entire functions (in fact polynomials) of $z$, and $E_J$ is
uniformly bounded for $0\le x\le a$ and $z$ in compact subsets of the plane.
Integration from zero to $a$ gives
\[
 \sum_{j=0}^{J-1}\frac{b_j(z)a^{j+1-z}}{j+1-z}
 +\int_0^a x^{J-z}E_J(x,z)\,dx.
\]
The second term is holomorphic for $\Re z<J+1$; the integral from $a$ to
one is entire.  Since $J$ is arbitrary, these formulas give the claimed
continuation.  At $z=k$ the residue is
$-b_{k-1}(k)=(-1)^k[u^{k-1}]\Gamma(1-u)^k$.
Finally
\[
 \log\Gamma(1-u)=\gamma u+\sum_{j\ge2}\frac{\zeta(j)}j u^j
\]
has strictly positive coefficients, so \eqref{eq:ck-lagrange} is positive.
\end{proof}

The pole at $k$ contributes $(-1)^{k-1}c_k/k^n$ to the Taylor coefficient
of $F$ at zero.  This explains the exponential scales in the known moment
asymptotics, but an infinite partial-fraction sum cannot be inferred
without controlling the growth of $c_k$.  That growth is decisive.

\subsection{Late coefficients and the inverse reciprocal gamma function}
Set $\phi(u)=\Gamma(1-u)$ and let $\tau$ be the unique point in $(0,1)$
satisfying
\begin{equation}
 -\tau\psi(1-\tau)=1.
 \label{eq:critical-tau}
\end{equation}
Equivalently $\psi(-\tau)=0$.  Define
\begin{equation}
 \rho=\frac{\tau}{\Gamma(1-\tau)},\qquad
 \alpha=-\log\rho,\qquad
 v=1+\tau^2\psi'(1-\tau),\qquad
 C=\frac{\tau}{\sqrt{2\pi v}}.
 \label{eq:critical-constants}
\end{equation}
Numerically,
\begin{align*}
 \tau&=0.50408300826445540925826930453330249895538518236858\ldots,\\
 \rho&=0.28211580900629804215379660062046372823230190757940\ldots,\\
 \alpha&=1.26543762211086561338203591116547399218544556554213\ldots,\\
 v&=2.27159994608893009911993485126846348661898954658806\ldots,\\
 C&=0.13342776130944528489602057053941782432703493711204\ldots.
\end{align*}
These displayed approximations are numerical evaluations; the theorem uses
the exact definitions \eqref{eq:critical-tau}--\eqref{eq:critical-constants}.

\begin{theorem}[Late-coefficient asymptotic]
\label{thm:late-coefficients}
The coefficients \eqref{eq:ck-lagrange} satisfy
\begin{equation}
 c_k=C\rho^{-k}k^{-3/2}\bigl(1+O(k^{-1})\bigr).
 \label{eq:ck-asymptotic}
\end{equation}
In fact they have a full expansion in inverse powers of $k$.
The power series $h(w)=\sum_{k\ge1}c_kw^k$ has radius of convergence
$\rho$, converges absolutely on $|w|\le\rho$, and satisfies
\[
 h(w)=w\Gamma(1-h(w)),\qquad h(\rho)=\tau.
\]
For every fixed real $n$, the formal moment series
$\sum_{k\ge1}(-1)^{k-1}c_k/k^n$ diverges, since its terms do not tend to zero.
\end{theorem}

\begin{proof}
Let $D=u\,d/du$.  The function
\[
 \mu(u)=D\log\phi(u)=-u\psi(1-u)
 =\gamma u+\sum_{j\ge2}\zeta(j)u^j
\]
is strictly increasing from zero to infinity on $(0,1)$.  Hence
\eqref{eq:critical-tau} has exactly one solution.
Moreover $\phi(u)>1$ there, so $0<\rho<1$ and $\alpha>0$.

Write $\phi(u)=\sum_{j\ge0}a_ju^j$ and introduce an integer-valued random
variable $X$ by
\[
 \Pr(X=j)=\frac{a_j\tau^j}{\phi(\tau)}.
\]
All its probabilities are positive, in particular those at $0$ and $1$;
it has lattice span one.  Analyticity of $\phi$ on $|u|<1$ gives an
exponential moment.  Its mean is $\mu(\tau)=1$ and its variance is
$D^2\log\phi(\tau)=1+\tau^2\psi'(1-\tau)=v>0$.
For independent copies and $S_k=X_1+\cdots+X_k$, coefficient extraction gives
\begin{equation}
 c_k=\frac{\tau\rho^{-k}}k\Pr(S_k=k-1).
 \label{eq:coefficient-probability}
\end{equation}
Fourier inversion gives
\[
 \Pr(S_k=k-1)=\frac1{2\pi}\int_{-\pi}^{\pi}
 \left(\frac{\phi(\tau e^{it})}{\phi(\tau)}e^{-it}\right)^k e^{it}\,dt.
\]
For $t\ne0$ modulo $2\pi$, the modulus of the bracket is strictly less
than one, because the distribution has positive mass at $0$ and $1$.
On a small interval around zero its logarithm is
$-vt^2/2+O(t^3)$.  Split off the complement of that interval, whose
contribution is exponentially small, and put $t=y/\sqrt{k}$ in the
remaining integral.  Taylor expansion, dominated by a Gaussian on a
shrinking central interval, gives
\[
 \Pr(S_k=k-1)=\frac1{\sqrt{2\pi vk}}\bigl(1+O(k^{-1})\bigr).
\]
For a fully quantitative split, use $|t|\le k^{-2/5}$, the band
$k^{-2/5}<|t|\le\delta$, and the fixed outer band.  On the middle
band the modulus is at most $e^{-c k^{1/5}}$; on the outer band it is
$O(q^k)$ for some $q<1$.  The nominal $k^{-1/2}$ correction on the
central band is odd and integrates to zero.  Taylor expansion to arbitrary
even order with its analytic remainder gives the full inverse-power expansion.  This elementary saddle argument proves
\eqref{eq:ck-asymptotic} and avoids any assumption about unexamined complex
critical values of the reciprocal gamma function.

Lagrange inversion gives the stated functional equation for $h$ near zero.
Equation \eqref{eq:ck-asymptotic} proves that its radius is $\rho$ and that
it converges absolutely at the boundary.  On $[0,\tau]$, the function
$f(u)=u/\phi(u)$ increases from zero to $\rho$, with positive derivative
before $\tau$, since $f'(u)=(1-\mu(u))/\phi(u)$.  Its inverse is the
positive real branch of $h$, hence continuity at the boundary gives
$h(\rho)=\tau$.  Finally \eqref{eq:ck-asymptotic} proves the stated divergence.
\end{proof}

For completeness, the first correction in \eqref{eq:ck-asymptotic} is
\[
 c_k=C\rho^{-k}k^{-3/2}\left(1+\frac{\beta_1}{k}+O(k^{-2})\right),
\]
where, with $\kappa_j=D^j\log\phi(\tau)$,
\begin{align*}
 \beta_1&=-\frac1{2v}+\frac{\kappa_3}{2v^2}
 +\frac{\kappa_4}{8v^2}-\frac{5\kappa_3^2}{24v^3}
 \\
 &=0.30196282878120524112458668271647811378573305248468\ldots.
\end{align*}
This follows by retaining the quartic and squared cubic terms, together
with the factor $e^{it}$, in the Fourier integral.

\subsection{An effective remainder with a growing number of terms}
\begin{theorem}[Explicit truncation bound]
\label{thm:effective-moment}
For integers $n\ge1$ and $K\ge1$ with $K\alpha<n$, put
\[
 R_{n,K}=\frac{M_n}{n!}-\sum_{k=1}^{K}\frac{(-1)^{k-1}c_k}{k^n}.
\]
Then, writing $\Gamma(n,z)=\int_z^\infty t^{n-1}e^{-t}\,dt$,
\begin{equation}
 |R_{n,K}|\le
 \frac{\alpha^n}{n!}\left(1+\frac{n\tau}{n-K\alpha}\right)
 +\frac{\tau\rho^{-(K+1)}}{(K+1)^n}
   \frac{\Gamma(n,(K+1)\alpha)}{\Gamma(n)}.
 \label{eq:remainder-exact}
\end{equation}
In particular, whenever $n\ge\max\{2,4\alpha^2\}$, the choice
$K=\lfloor n/\alpha-\sqrt n\rfloor$ is positive and gives
\begin{equation}
 |R_{n,K}|\le
 \left(\frac1{\sqrt{2\pi n}}+
       \frac{\tau}{\alpha\sqrt{2\pi}}+
       \tau e^{2\alpha^2}\right)
 \left(\frac{e\alpha}{n}\right)^n.
 \label{eq:remainder-optimized}
\end{equation}
\end{theorem}
In particular every integer $n\ge7$ meets the displayed hypothesis.
Indeed $\rho\ge1/(2\sqrt\pi)$ by evaluating $u/\Gamma(1-u)$
at $u=1/2$, and therefore $\alpha\le\log(2\sqrt\pi)<13/10$,
so $4\alpha^2<169/25<7$.

\begin{proof}
The function $\log\Gamma(x)$ decreases strictly from infinity to zero on
$(0,1]$, because $\psi'(x)>0$ and $\psi(1)=-\gamma<0$.
Let $x(y)\in(0,1]$ be its inverse.  The layer-cake formula, or integration
by parts after this monotone change of variables, gives
\begin{equation}
 \frac{M_n}{n!}=\frac1{\Gamma(n)}\int_0^\infty y^{n-1}x(y)\,dy.
 \label{eq:moment-inverse}
\end{equation}
For $y\ge\alpha$, Lagrange inversion and continuity on the convergence
circle give
\[
 x(y)=-h(-e^{-y})=\sum_{k\ge1}(-1)^{k-1}c_ke^{-ky}.
\]
The identity initially holds for sufficiently large $y$ and extends along
the negative real inverse branch.  This branch cannot fail before
$w=-\rho$: for $0<x\le1$, the derivative of $1/\Gamma(x)$ is
$-\psi(x)/\Gamma(x)>0$.

For $y\ge\alpha$, positivity and $\sum c_k\rho^k=\tau$ give
\[
 \left|x(y)-\sum_{k=1}^{K}(-1)^{k-1}c_ke^{-ky}\right|
 \le\tau\rho^{-(K+1)}e^{-(K+1)y}.
\]
Integrating this on $[\alpha,\infty)$ produces the second term of
\eqref{eq:remainder-exact}.
For the interval $[0,\alpha]$ use $0<x(y)\le1$ and, since $n>k\alpha$,
\[
 \int_0^\alpha y^{n-1}e^{-ky}\,dy
 \le\frac{\alpha^n e^{-k\alpha}}{n-k\alpha}.
\]
To see this inequality, differentiate $y^ne^{-ky}$ and use
$n-ky\ge n-k\alpha$.  Sum over $1\le k\le K$ and bound
$\sum c_ke^{-k\alpha}\le\tau$ to obtain the first term.

For the stated choice of $K$, $n-K\alpha\ge\alpha\sqrt n$.
Let $\theta=(K+1)\alpha/n$.  Then
$1-\alpha/\sqrt n<\theta\le1$ and the hypothesis ensures $\theta\ge1/2$.
Since
\[
 \theta-1-\log\theta\le2(1-\theta)^2\qquad(1/2\le\theta\le1),
\]
we have
\[
 \rho^{-(K+1)}(K+1)^{-n}
 =\left(\frac{e\alpha}{n}\right)^n
   e^{n(\theta-1-\log\theta)}
 \le e^{2\alpha^2}\left(\frac{e\alpha}{n}\right)^n.
\]
Finally use $\Gamma(n,z)/\Gamma(n)\le1$ and the lower Stirling bound
$n!\ge\sqrt{2\pi n}(n/e)^n$.
\end{proof}

For each fixed $K$, retaining one additional term in
\eqref{eq:remainder-exact} recovers the usual asymptotic interpretation.
The stronger content of \eqref{eq:remainder-optimized} is that $K$ may grow
with $n$.  The bound has the scale $(e\alpha/n)^n$.
The late-coefficient theorem puts the least term near
$(n+3/2)/\alpha$ and gives its magnitude
\[
 C\left(\frac{e\alpha}{n+3/2}\right)^{n+3/2}(1+O(n^{-1})).
\]
The rigorously proved remainder bound is within a polynomial factor of
this scale.  A least-term estimate alone would not prove a remainder bound;
the inverse-function integral is what supplies that missing control.

The following values are numerical diagnostics from an independent integral
in the variable $t=-\log x$, computed at $180$ decimal digits.  The
rightmost column evaluates the explicit bound in
\eqref{eq:remainder-optimized}; the inequality itself is proved above.
\begin{center}
\begin{tabular}{rrrr}
\hline
$n$ & $K$ & Observed $|R_{n,K}|$ & Explicit upper bound \\
\hline
$8$ & $3$ & $3.1162\cdot10^{-5}$ & $1.4837\cdot10^{-2}$ \\
$20$ & $11$ & $2.1101\cdot10^{-18}$ & $6.4881\cdot10^{-15}$ \\
$40$ & $25$ & $2.8961\cdot10^{-46}$ & $3.0208\cdot10^{-42}$ \\
$80$ & $54$ & $1.8102\cdot10^{-113}$ & $5.9717\cdot10^{-109}$ \\
\hline
\end{tabular}
\end{center}
The displayed bound decimals have been rounded upward.  They serve as a
scale comparison; exact theorem constants are defined by
\eqref{eq:critical-constants}.  Full numerical data are in
\texttt{moment\_diagnostics.json}.

\subsection{A 112-digit rational certificate for the cubic moment}
The file \texttt{certify\_cubic.py} implements a separate high-precision
calculation using integer outward intervals.  No floating-point special
function is trusted.  Put $a=1/2$, $\ell=\log2$ and
\[
 b_1=\gamma,\qquad b_j=\frac{\zeta(j)}j\ (j\ge2),\qquad
 B_M(x)=\sum_{j=1}^{M}b_jx^j.
\]
On $0<x\le a$ we approximate
$\log\Gamma(x)$ by $-\log x+B_M(-x)$; on the upper half of the interval,
put $t=1-x$ and approximate $\log\Gamma(1-t)$ by $B_M(t)$.
Both errors are at most
\[
 D x^{M+1},\qquad D=\frac{2\zeta(2)}{M+1},
\]
with the corresponding variable $x$ or $t$.  This follows from the Taylor
series and $\zeta(j)\le\zeta(2)$ for $j\ge2$.
Define the exact elementary integrals
\[
 J_{j,p}=\int_0^{1/2}x^j(-\log x)^p dx
 =\frac{2^{-j-1}p!}{(j+1)^{p+1}}
 \sum_{r=0}^{p}\frac{((j+1)\ell)^r}{r!}.
\]
If $b^{*r}_j$ are the coefficients of $B_M^r$, the integrated polynomial
approximation is
\[
 V_M=J_{0,3}+3\sum_{j=1}^{M}(-1)^jb_jJ_{j,2}
 +3\sum_{j=2}^{2M}(-1)^jb^{*2}_jJ_{j,1}
 +2\sum_{\substack{j\le3M\\j\text{ even}}}b^{*3}_jJ_{j,0}.
\]
The last sum starts at $j=4$.  Since $0\le\log\Gamma(x)\le-\log x$
on the lower half and $0\le\log\Gamma(1-t)\le1$ on the upper half,
\begin{align}
 |M_3-V_M|\le E_M:={}&3D(J_{M+1,2}+J_{M+1,0})\notag\\
 &+3D^2(J_{2M+2,1}+J_{2M+2,0})+2D^3J_{3M+3,0}.
 \label{eq:cubic-certified-error}
\end{align}
Here $\Gamma(1+x)\le1$ follows from log-convexity between $1$ and $2$;
the upper-half bound follows from $\gamma<1$, $\zeta(j)<2$, and
$\sum_{j\ge2}(1/2)^j/j\le1/4$.

For reproducibility, the constant intervals come from the positive
atanh series for $\log2$, the Euler--Maclaurin formula for $\zeta(s)$ at
$N=128$, and the harmonic-number formula for $\gamma$.  With $m=96$, their
respective analytic error bounds are
\[
 \frac{|B_{2m}|}{(2m)!}(s)_{2m-1}N^{-s-2m+1},\qquad
 \frac{|B_{2m}|}{2mN^{2m}}.
\]
The sharper first-omitted-term zeta bound uses the Stieltjes remainder
property, not just an absolute periodic-Bernoulli estimate.  Namely,
\[
 \frac1{1-e^{-t}}=\frac1t+\frac12+2t\sum_{r\ge1}
 \frac1{t^2+4\pi^2r^2}.
\]
Expanding each denominator geometrically through the term preceding
$B_{2m}$ gives a signed remainder bounded in absolute value by
$|B_{2m}|t^{2m-1}/(2m)!$.  Integrate this against
$t^{s-1}e^{-Nt}/\Gamma(s)$ to obtain the stated bound.  This argument
also specifies exactly which Euler--Maclaurin terms are retained:
$j=1,\ldots,m-1$.
The gamma bound follows by expanding the geometric factor in Binet's
integral for the digamma function.  For $s\ge64$ the code instead brackets
the tail of the direct positive sum by the two adjacent integrals.
All rational operations round outwards at a fixed denominator $10^{130}$.
Exact Bernoulli numbers are provided by SymPy; no numerical Bernoulli
approximation is involved.

At $M=360$, the total interval has width less than $2.387\cdot10^{-113}$.
In particular it certifies the displayed digits
\[
 M_3=5.74038880722947428001957168810246146296101300745487\ldots.
\]
The full lower and upper rational endpoints, and the analytic error bound,
are stored in \texttt{cubic\_certificate.json}.  An independent
150-digit mpmath quadrature lies inside this proved interval; that numerical
check is a diagnostic, whereas \eqref{eq:cubic-certified-error} and the
outward arithmetic supply the certificate.

\section{Reproducible computation and the meaning of a certificate}
\label{sec:computation}

The archive is organized so that the analytic argument can be read
without executing code, while each finite algebraic step can be replayed
without trusting a decimal fit. There are three distinct kinds of
computation. Exact rational word algebra verifies a shuffle certificate.
Outward interval arithmetic, together with a proved analytic tail bound,
certifies a numerical enclosure. Ordinary high-precision quadrature checks
an implementation and its conventions. These roles should remain
distinguishable when the files are integrated into ProveIt.

\subsection{Exact word algebra and algebraic substitutions}

The program \code{code/shuffle_certificates.py} represents a binary word
as a tuple and a shuffle as a dictionary of integer multiplicities.
It recursively shuffles two words, computes the nonincreasing Lyndon
factorization, and solves the resulting triangular equation using
rational numbers. The decisive replay is to expand the resulting
polynomial back into words. A successful replay returns exactly the
original word with coefficient one and no other terms.

At weights one through nine, this procedure re-expands all \(511\)
nonempty words ending in one. It also enumerates all \(45\) nonempty
bigraded sectors and checks their Lyndon counts against
Theorem~\ref{thm:witt}. The recorded data contain every polynomial normal
form of the ten weight-six triples, all four weight-five doubles, and
the two-product certificate of Corollary~\ref{cor:wt5}.
This exhaustive finite check supports the implementation; the
all-weight statement rests on the proof of the theorem.

The depth verifier implements the binomial formula independently of the
target expressions. Exact symbolic subtraction confirms every coefficient
in the five weight-six and three weight-four source formulas.
It also emits the parity-reducible Gaussian component of every double
through weight twelve: imaginary parts at even weight and real parts at
odd weight. Thus the archive includes further explicit identities as
reusable data, with no integer-relation recognition step.

The ladder verifier works in
\(\Q[r]/(r^3+r-1)\). Rational functions are reduced using the inverse
of their denominator modulo the defining polynomial. Its \(35\) exact
checks cover the six factorizations, the eighteen Kummer arguments,
the eight Rogers arguments, the rational row operation, and both
logarithmic-tail cancellations. In particular, the finite certificate
does not attempt to evaluate a transcendental function in the number
field: it verifies the arguments and coefficients of functional equations
whose analytic validity is stated in the article.

\subsection{Independent numerical diagnostics}

For the same-argument double family, a useful independent evaluator is
\begin{equation}
F_{a,b}(z)=\frac{z}{(a-1)!}
\int_0^1\frac{(-\log t)^{a-1}\Li_b(zt)}{1-zt}\,dt,
\qquad |z|=1,\ z\ne1.
\label{eq:double-quadrature}
\end{equation}
It follows first by absolutely convergent series expansion for \(|z|<1\),
then by the boundary limit. This evaluator uses neither the binomial
reduction nor a fitted constant basket.
For the inverse-argument mixed family the corresponding integral replaces
\(\Li_b(zt)\) by \(\Li_b(t)\).
For the arbitrary-depth one-two family the verifier uses the integral
with the zero letter held fixed, before expanding its outer factor.

The resulting diagnostics comprise \(64\) checks at \(65\) decimal
working digits for the same-argument double and one-two formulas,
including both Gaussian and Eisenstein points and the non-cyclotomic
point \(e^{i\sqrt2}\). Their largest absolute residual was
\(1.081\cdot10^{-64}\) when rounded upward.
A separate inherited-mixed audit contains eighteen general-formula
quadratures and four explicit specialization checks, with maximum
residual less than \(1.365\cdot10^{-65}\).
The cubic and complementary-power ladders were evaluated at \(100\)
and \(200\) working digits, including comparison with directly summed
convergent classical-polylogarithm series. The cubic residual at
\(200\) digits was below \(2.62\cdot10^{-200}\); the independent
series comparison differed by less than \(9.91\cdot10^{-191}\).

For the remaining \(S_4\) conjecture, an independent real integral is
\[
S_4=-\frac16\int_0^1
\frac{(-\log t)^3\log(1+t^2)}{1+t^2}\,dt.
\]
It gives
\[
S_4=-0.0104934562973350434533567510062017519635106188773\ldots.
\]
At \(50\) working digits the discrepancy from
\eqref{eq:S4-short} was below \(5.0\cdot10^{-51}\).
This is evidence for the conjecture, and is deliberately excluded from
the proved-identity count. The archive includes a small independent
script for this comparison.

None of these residuals is an interval-certified error bound. In
particular, a residual at approximately the working precision is
consistent with cancellation and finite-precision rounding.
The independent formulas and the analytic proofs, rather than the
number of matching digits, determine the mathematical status.

\subsection{The interval certificate and the asymptotic diagnostics}

The cubic moment calculation has a stronger status. All endpoints in
\code{code/moments/cubic_certificate.json} are exact rational numbers
with denominator \(10^{130}\), displayed in decimal form. The
certificate uses no floating-point special-function values. The interval
width is less than \(2.387\cdot10^{-113}\), and its endpoints share
\(112\) digits after the decimal point. This certifies the common
truncated decimal prefix; a claim about rounded digits would instead
require checking the relevant rounding boundary.

The certificate was replayed independently and reproduced every recorded
field exactly. Separate quadratures at \(150\) and \(180\) working
digits lie inside the interval. The quadratures provide a useful
cross-check but are not used to construct the rational enclosure.

The coefficient asymptotic and moment bounds were also compared with
direct evaluations. Figure~\ref{fig:moment-asymptotics} shows two
features that matter for using the expansion. The coefficient
approximation improves in the predicted inverse powers of \(k\).
Nevertheless, at each fixed \(n\), the terms \(c_k/k^n\) eventually
increase. It is therefore essential to specify a truncation rule,
such as the one in Theorem~\ref{thm:effective-moment}.

\begin{figure}[htbp]
\centering
\includegraphics[width=\linewidth]{figures/moment_asymptotics.pdf}
\caption{Numerical diagnostics for the proved coefficient and moment
asymptotics. Left: absolute relative errors of
\(C\rho^{-k}k^{-3/2}\) and
\(C\rho^{-k}k^{-3/2}(1+\beta_1/k)\), using coefficients computed from
their positive defining recurrence. Right:
\(\log_{10}(c_k/k^n)\) for \(n=20,40,80\). Dotted vertical lines
mark \(k=n/\alpha\); circles mark the sampled minima.
All plotted quantities are dimensionless. The figure illustrates
the theorems and supplies no independent error certificate.}
\label{fig:moment-asymptotics}
\end{figure}

\subsection{Archive structure and reproduction}

The top-level \code{README.md} gives installation, build, verification,
and integration instructions. The article's master source is
\code{article/polylogarithms_exact_reductions.tex}; all included
sections and the vector figure are supplied. The verification programs
are grouped under \code{code/}, with their recorded outputs alongside
the relevant program or under \code{data/}. Source provenance includes
the pinned commit, content hashes, and a status register mapping each
proposed manuscript update to its source equation or section.

From the archive root, the essential commands are
\begin{quote}\small\ttfamily
python -m pip install -r requirements.txt\\
python code/run\_checks.py --fast\\
python code/run\_checks.py --full\\
make pdf
\end{quote}
The fast command replays exact algebra and the rational interval
certificate. The full command additionally runs the independent
quadratures and regenerates the diagnostic figure. It may take several
minutes; no network access is required after installing the dependencies.
The complete recorded diagnostics remain available without rerunning
them. The PDF is built with ordinary \LaTeX{} packages and contains
no external data dependency.

The source manifest is a checkpoint, not an assertion that later
repository revisions have the same status. Integration should preserve
the proved/conjectured distinction, the descending-index convention,
the branch choices, and the qualification that shuffle quotient
dimensions are formal. The supplied register is intended to make
that editorial task precise.

\section{Further research: specific questions and workable next steps}
\label{sec:research}

The results change the most useful next questions. Repeating searches
for the eleven settled identities would add little. The next stage
should seek additional functional relations, effective bounds, and
well-specified arithmetic statements. The following problems distinguish
a plausible method from the mathematical conclusion it would need
to establish.

\subsection{First priority: the remaining Gaussian bridge}

\paragraph{1. Prove the \(S_4\) relation with an exact colored certificate.}
After the two shuffle eliminations, Conjecture~\ref{conj:S4} is the
single concrete target
\[
7\Im\Li_{4,1}(i,-1)
=51g_{41}+24g_{32}+\frac{19\pi^5}{512}
-14\beta(4)\log2.
\]
This displayed equation is still conjectural. Its second argument
is \(-1\), so it is not one of the already reduced inverse-argument
values \(\Li_{a,b}(i,-i)\).
A promising approach is to work with iterated integrals on the alphabet
\(\{0,1,-1,i,-i\}\), include explicitly regularized shuffle and stuffle
relations, and track the additional transformations of the punctured
line. The desired output is a rational linear combination of justified
relations with all boundary constants recorded.
An exact certificate would be a substantial improvement over another
decimal match.

\paragraph{2. Audit the colored relation system before enlarging the
constant basket.}
The inherited mixed formulas have modest integer coefficients, yet the
current manuscript reports them as apparent non-reductions.
Reproduce that computation from its value-generation stage, keeping
argument ordering and the parity component explicit.
At each weight record the exact alphabet, admitted product rows,
regularization prescriptions, distribution identities, and the
coefficient-height limit used by any numerical search.
Then compare the resulting exact row space with the explicit formulas
in Section~\ref{sec:mixed-reconciliation}.
This should determine whether the discrepancy comes from the numerical
values, a missing basis element, or an incomplete relation system.
An unsuccessful finite search should never be promoted to a statement
of period independence.

\paragraph{3. Determine what survives at weight six and depth three.}
The classes of \(F_{4,1,1},F_{3,2,1},F_{3,1,2}\) are an exact basis
of the formal one-variable shuffle quotient.
At \(i\), additional colored functional equations can reduce their
span further. The next question is the rank after specifying those
additional relations and the allowed lower-depth algebra.
The all-depth one-two theorem does not answer this: these indices
contain three zero letters in their words.
A useful outcome would be either explicit reductions or a certified
presentation of a larger, precisely defined formal quotient.
Neither outcome by itself proves that the numerical values are
linearly independent.

\subsection{Algebraic ladders beyond the two local proofs}

\paragraph{4. Produce local proofs of the higher golden ladders.}
The manuscript's higher golden formulas should receive the same kind
of argument/substitution and constant-term certificate as the cubic
identity. Start with the lowest unresolved weight and its full
functional-equation closure.
A vanishing symbol or differential can narrow the problem, but the
remaining products, zeta terms, and endpoint constants still require
determination. In particular, an identity valid only at an isolated
algebraic number cannot simply be integrated in that number to obtain
a higher-weight identity. A valid deformation must preserve the
functional relations being used.
Success here means a branch-consistent analytic proof of one exact
recorded formula, not only a lower-height numerical relation.

\paragraph{5. Classify useful complementary-power specializations.}
Theorem~\ref{thm:complementary-powers} starts from any
\(\rho^a+\rho^b=1\), but different exponent pairs may describe the
same field or the same base after a power substitution.
For bounded degree and exponent height, classify these redundancies
and determine which Kummer substitutions close on a small set of
powers of \(\rho\).
The cubic proof suggests looking for a closure that introduces a few
auxiliary powers and then cancels them through short rational rows.
The output can be an exact finite database of argument identities
over number fields, together with the corresponding polylogarithm
certificates. This is a constructive classification problem even when
period independence is out of reach.

\paragraph{6. Develop a reusable algebraic functional-equation verifier.}
Generalize the current quotient-ring calculation to arbitrary
irreducible polynomials, with an isolating interval or another exact
embedding specification. The verifier should check denominator
nonvanishing, identify the real or complex embedding, and retain
sign information needed for ordinary logarithms.
For single-valued functions, argument reduction is often sufficient;
for ordinary polylogarithms, continuation paths and inversion
corrections must remain part of the certificate.
The natural first benchmark is to reproduce the cubic certificate
without any special-case algebra in the program.

\subsection{Quantitative analysis of log-gamma moments}

\paragraph{7. Find the signed remainder near the least term.}
Theorem~\ref{thm:effective-moment} gives an unsigned bound of scale
\((e\alpha/n)^n\), within a polynomial factor of the smallest term.
It does not determine the leading signed error or prove an optimal
constant. Study \(R_{n,K}\) for
\(K=(n+3/2)/\alpha+O(1)\), allowing the integer rounding and the
alternation in \(k\) to enter explicitly.
The exact inverse-function integral is a starting point; the current
positive majorant discards cancellations that may control a sharper
answer. Any proposed asymptotic should first be checked over both
parities of \(K\), then proved uniformly in the rounding parameter.

\paragraph{8. Extend the theorem to weighted and reflected moments.}
For \(\beta>-1\), consider
\[
M_n(\beta)=\int_0^1x^\beta\log^n\Gamma(x)\,dx.
\]
The same layer-cake argument gives the exact starting identity
\[
\frac{M_n(\beta)}{n!}
=\frac{1}{(\beta+1)\Gamma(n)}
\int_0^\infty y^{n-1}x(y)^{\beta+1}\,dy,\qquad n\ge1.
\]
This identity follows by integrating \(x^\beta\) over the sublevel
interval \(0<x<x(y)\); it is not a conjecture.
For integer \(\beta+1\), Lagrange--B\"urmann coefficients of powers
of \(h\) suggest a direct extension of the coefficient analysis.
For noninteger powers, fractional leading exponents and branch
normalizations must be handled explicitly.
Mixed reflected moments involving
\(\log\Gamma(x)\) and \(\log\Gamma(1-x)\) add a second endpoint
mechanism and should be treated separately.

\paragraph{9. Determine the next inverse-gamma singularities.}
The local-limit proof finds the dominant growth without assuming a
global inverse-branch picture. A further analytic study should determine
the critical values of the reciprocal-gamma map accessible from this
branch, their local types, and the continuation paths that reach them.
That information could organize subsequent exponentially small
corrections and explain the behavior of higher-order truncation.
Computing roots of the digamma function numerically is a useful
exploration, but is not a proof of completeness or of which critical
value controls a given branch.

\paragraph{10. Reduce the known Tornheim coordinates in a specified
comparison algebra.}
The existence of a named Tornheim representation for \(M_3\) is settled
by \cite{BBB}. A meaningful remaining problem is, for example, whether
the particular combination \(2T_{13}-T_{12}\) at \((1,1,1)\) admits
a reduction using an explicitly stated set of zeta derivatives,
logarithms, and other standard constants.
The set must be fixed before a numerical search, and failure to
find a relation must retain its precision and height qualifications.
Independently, improve the triangular tail estimate
\eqref{eq:tornheim-tail} by summing in \(q=m+n\) and applying
an Euler--Maclaurin expansion to the resulting harmonic-logarithmic
coefficients. This could certify individual Tornheim combinations
without using the already evaluated cubic integral.

\subsection{A longer-term exactness programme}

\paragraph{11. Separate functional reduction from arithmetic
independence.}
The present results supply upper bounds and explicit reductions.
Lower bounds for numerical period spaces require different input.
One direction is to formulate the same questions in a motivic or
other formal period algebra, where a structural dimension can
sometimes be established, and then state exactly what comparison
or period conjecture would transfer that dimension to numbers.
The formal shuffle quotient of Theorem~\ref{thm:witt} is a useful
combinatorial baseline, but it contains too few relations to serve
as that arithmetic model on its own.

\paragraph{12. Move the smallest certificates into a proof assistant.}
The first targets should be the finite word-shuffle normal forms,
the binomial matrix inverse, and the rational algebraic substitutions.
They need only combinatorics and exact arithmetic.
The next layer consists of analytic lemmas: convergence of the
boundary Fourier convolution, the stated branches of polylogarithm
inversion, and a rigorous positive-tail implementation for the cubic
moment. A machine-checked rational row is valuable only when the
functional equations it combines have been justified in the same
convention. This staged approach keeps the trusted analytic base
visible and makes each new certificate reusable.

The common objective is a research record in which a numerical
discovery can be followed by a short exact witness, and in which the
questions that remain are stated at the level actually supported
by the evidence. The two-product shuffle, the algebraic ladder rows,
and the inverse-gamma remainder estimate provide three concrete
models for that next stage.

\clearpage
\begin{thebibliography}{99}
\addcontentsline{toc}{section}{References}

\bibitem{manuscript}
V.~Reshetnikov, \emph{Polylogarithms and their Arithmetic Bridges},
ProveIt unified manuscript, source checkpoint
\code{afed07429d3d37eceb6c8e9e54cf4da2d3f39d53}, 9 October 2026.
\url{https://github.com/VladimirReshetnikov/ProveIt/tree/afed07429d3d37eceb6c8e9e54cf4da2d3f39d53/Analysis/Polylogarithms/docs/manuscript}.

\bibitem{Radford}
D.~E.~Radford, A natural ring basis for the shuffle algebra and an
application to group schemes, \emph{Journal of Algebra} \textbf{58}
(1979), no.~2, 432--454.
\href{https://doi.org/10.1016/0021-8693(79)90171-6}{doi:10.1016/0021-8693(79)90171-6}.

\bibitem{Panzer}
E.~Panzer, The parity theorem for multiple polylogarithms,
\emph{Journal of Number Theory} \textbf{172} (2017), 93--113.
\href{https://doi.org/10.1016/j.jnt.2016.08.004}{doi:10.1016/j.jnt.2016.08.004}.
Author preprint: \url{https://arxiv.org/abs/1512.04482}.

\bibitem{Umezawa}
R.~Umezawa, \emph{An explicit parity theorem for multiple polylogarithms},
arXiv:2508.02040, 2025.
\url{https://arxiv.org/abs/2508.02040}.

\bibitem{Zagier}
D.~Zagier, Special values and functional equations of polylogarithms,
in L.~Lewin (ed.), \emph{Structural Properties of Polylogarithms},
Mathematical Surveys and Monographs \textbf{37}, American Mathematical
Society, Providence, RI, 1991, 377--400.
Author copy:
\url{https://people.mpim-bonn.mpg.de/zagier/files/tex/LewinPolylogarithms/fulltext.pdf}.

\bibitem{NakamuraShiraishi}
H.~Nakamura and D.~Shiraishi, Landen's trilogarithm functional equation
and \(\ell\)-adic Galois multiple polylogarithms,
in M.~Morishita, H.~Nakamura, and J.~Ueki (eds.),
\emph{Low Dimensional Topology and Number Theory},
Springer Proceedings in Mathematics \& Statistics \textbf{456},
Springer, Singapore, 2025, 237--262.
\href{https://doi.org/10.1007/978-981-97-3778-9_8}{doi:10.1007/978-981-97-3778-9\_8}.
Author preprint: \url{https://arxiv.org/abs/2210.17182}.

\bibitem{ACEMKM}
T.~Amdeberhan, M.~W.~Coffey, O.~Espinosa, C.~Koutschan,
D.~V.~Manna, and V.~H.~Moll, Integrals of powers of loggamma,
\emph{Proceedings of the American Mathematical Society}
\textbf{139} (2011), no.~2, 535--545.
\href{https://doi.org/10.1090/S0002-9939-2010-10589-0}{doi:10.1090/S0002-9939-2010-10589-0}.
Author preprint:
\url{https://www.math.tulane.edu/~vhm/papers_html/lg-subm.pdf}.

\bibitem{BBB}
D.~H.~Bailey, D.~Borwein, and J.~M.~Borwein,
On Eulerian log-gamma integrals and Tornheim--Witten zeta functions,
\emph{The Ramanujan Journal} \textbf{36} (2015), nos.~1--2, 43--68.
\href{https://doi.org/10.1007/s11139-012-9427-1}{doi:10.1007/s11139-012-9427-1}.
Author preprint:
\url{https://www.davidhbailey.com/dhbpapers/log-gamma.pdf}.

\end{thebibliography}

\end{document}
